\documentclass[landscape]{article}

% Lane comparison for the italic base-anchor candidate X (2026-11-10).
% Tables: .venv/bin/python run_print_lane_comparison.py (from the repo root).
% PDF: cd tex && xelatex lane-comparison.tex

\usepackage[table]{xcolor}
\usepackage{fontspec}
\usepackage{geometry}
\usepackage{caption}
\usepackage{array}

\geometry{a4paper, left=1.5cm, right=1.5cm, top=1.5cm, bottom=1.5cm}
\setlength{\parindent}{0pt}
\addtolength{\parskip}{2pt}

% always use `Path = ../fonts/`, always include the extension
\setmainfont{Libertinus Serif}[
	Path = ../fonts/,
	UprightFont    = LibertinusSerif-Regular.otf,
	ItalicFont     = LibertinusSerif-Italic.otf,
	BoldFont       = LibertinusSerif-Semibold.otf,
	BoldItalicFont = LibertinusSerif-SemiboldItalic.otf
]

\newfontfamily\LibertinusSerifPatch[
	Path = ../fonts/,
	UprightFont    = LibertinusSerif-Regular-patch.otf,
	ItalicFont     = LibertinusSerif-Italic-patch.otf,
	BoldFont       = LibertinusSerif-Semibold-patch.otf,
	BoldItalicFont = LibertinusSerif-SemiboldItalic-patch.otf
]{Libertinus Serif Patch}

	% combining grapheme joiner
\newcommand{\cgj}{\char"034F }

\begin{document}

\section*{Italic lanes for the base-anchor candidate $X$}

A \emph{lane} is a strip at the italic angle $\theta$ (12° italic, 11.5° semibold italic) with a left and a right edge.
A line at the italic angle through a point $(x, y)$ meets the baseline at $x - y\tan\theta$,
so a lane is given by where its edges meet the baseline, and its midline at height $Y$ is
the midpoint of those two values plus $Y\tan\theta$.
For a glyph with outline $\mathcal{O}$ and advance width $W$:

\begin{tabular}{lllll}
lane & left edge & right edge & midline at $Y$ & \\
\hline
advance lane & $0$ & $W$ & $W/2 + Y\tan\theta$ & the model \\
ink lane & $L = \min_{\mathcal{O}} (x - y\tan\theta)$ & $R = \max_{\mathcal{O}} (x - y\tan\theta)$ & $C + Y\tan\theta$, \ $C = (L+R)/2$ & the mark-anchor model's axis \\
ideal lane & $L^*$ & $R^*$ & $(L^*+R^*)/2 + Y\tan\theta$ & where the glyph lives; not computed \\
\end{tabular}

The ink lane is the narrowest strip at the italic angle that contains the outline:
its edges are lines at the italic angle that just touch the outline, as $X_\mathrm{min}$ and $X_\mathrm{max}$ are upright lines that just touch it.
With $\theta = 0$, $L = X_\mathrm{min}$ and $R = X_\mathrm{max}$, so in roman both lanes give the bbox midpoint, the roman candidate; only the italic styles are compared.
Italic n: $W = 518$, $L = 17$ (the entry serif at the top left), $R = 502$ (the outstroke at the bottom right), so $C = 259 = W/2$ and the lanes agree.
They differ where the ink sits off-centre in the advance width.

The candidate is $X = $ midline at the candidate's $Y$, minus $K$ (semibold italic: $K = 150$ above, $-12$ below; italic: $0$).
$Y$ is the same for both lanes.

\subsection*{Summary}

$\delta X$ = anchor $X$ minus candidate $X$.
\emph{designer}: the anchors of the original fonts (the fontmetrics); they were probably placed by advance width, so they favour the advance lane.
\emph{no derived or precomposed}: without the letters of \verb|data/derived_letters.py| (which inherit from their base letter) and canonically decomposed letters.
\emph{curated}: the human-curated base anchors (\verb|data/fontanchors_human/|, \verb|bases| by codepoint), with the candidate from the original glyph.
Columns per lane: median and mean $|\delta X|$, how many $|\delta X| \le 10$, median signed $\delta X$.
\emph{closer}: how many anchors each lane's candidate is closer to (rounded).

{\footnotesize
\begin{tabular}{lllrrrrrrrrrrrr}
style & anchors & class & $n$ & \multicolumn{4}{c}{advance lane} & \multicolumn{4}{c}{ink lane} & \multicolumn{3}{c}{closer} \\
 & & & & med & mean & ${\le}10$ & sgn & med & mean & ${\le}10$ & sgn & adv & ink & tie \\
\hline
it & designer, all & a & 82 & 26.4 & 69.4 & 21 & -26 & 27.2 & 69.9 & 22 & -25 & 41 & 37 & 4 \\
it & designer, all & b & 108 & 18.6 & 31.4 & 30 & 14 & 22.5 & 34.6 & 23 & 13 & 62 & 42 & 4 \\
it & designer, no derived or precomposed & a & 73 & 21.6 & 56.5 & 21 & -22 & 23.7 & 56.0 & 22 & -19 & 36 & 33 & 4 \\
it & designer, no derived or precomposed & b & 91 & 17.4 & 30.6 & 26 & 12 & 21.8 & 32.4 & 21 & 11 & 49 & 38 & 4 \\
it & curated & a & 30 & 20.0 & 35.9 & 5 & 1 & 22.5 & 40.0 & 8 & -3 & 16 & 13 & 1 \\
it & curated & b & 22 & 24.4 & 37.3 & 5 & -8 & 24.5 & 41.3 & 5 & -10 & 14 & 7 & 1 \\
\hline
si & designer, all & a & 91 & 21.9 & 34.0 & 17 & -10 & 24.9 & 36.1 & 19 & -8 & 54 & 36 & 1 \\
si & designer, all & b & 109 & 18.1 & 25.8 & 29 & -3 & 22.6 & 29.4 & 26 & -5 & 71 & 35 & 3 \\
si & designer, no derived or precomposed & a & 79 & 21.6 & 33.0 & 14 & -7 & 24.8 & 33.9 & 16 & -8 & 45 & 33 & 1 \\
si & designer, no derived or precomposed & b & 92 & 17.9 & 25.5 & 24 & -4 & 22.0 & 27.6 & 21 & -5 & 59 & 30 & 3 \\
si & curated & a & 22 & 16.7 & 34.8 & 7 & 1 & 18.6 & 40.4 & 4 & -2 & 13 & 9 & 0 \\
si & curated & b & 19 & 27.8 & 37.3 & 5 & -15 & 22.6 & 37.8 & 4 & -20 & 9 & 10 & 0 \\
\hline
\end{tabular}

}

\subsection*{Table rows}

hex, cat: codepoint and Unicode category.
it, si: the base alone, with dot above, and with macron below (via CGJ), original font; patch: the same in the patched font.
a, b: the designer's above and below anchor $X, Y$.
a?~adv, b?~adv: candidate from the advance lane; a$\delta$~adv, b$\delta$~adv: anchor minus that candidate.
a?~ink, b?~ink and a$\delta$~ink, b$\delta$~ink: the same for the ink lane.
Of the two $\delta$ cells, the one with the smaller $|\delta X|$ is bold.

\clearpage


\begin{table}[p]
\captionsetup{justification=raggedright, singlelinecheck=false}
\caption{Lane comparison for A--Z 1/2}
\label{table:lanes_AZ_1}
\setlength{\tabcolsep}{3.5pt}

{\footnotesize
\begin{tabular}{lrrrrrrrrrrrrr}
hex & 0041 & 0042 & 0043 & 0044 & 0045 & 0046 & 0047 & 0048 & 0049 & 004A & 004B & 004C & 004D \\
cat & Lu & Lu & Lu & Lu & Lu & Lu & Lu & Lu & Lu & Lu & Lu & Lu & Lu \\
\hline it\rule[-0.5em]{0pt}{1.9em} & {\itshape \char"0041 \char"0041\cgj\char"0307 \char"0041\cgj\char"0331} & {\itshape \char"0042 \char"0042\cgj\char"0307 \char"0042\cgj\char"0331} & {\itshape \char"0043 \char"0043\cgj\char"0307 \char"0043\cgj\char"0331} & {\itshape \char"0044 \char"0044\cgj\char"0307 \char"0044\cgj\char"0331} & {\itshape \char"0045 \char"0045\cgj\char"0307 \char"0045\cgj\char"0331} & {\itshape \char"0046 \char"0046\cgj\char"0307 \char"0046\cgj\char"0331} & {\itshape \char"0047 \char"0047\cgj\char"0307 \char"0047\cgj\char"0331} & {\itshape \char"0048 \char"0048\cgj\char"0307 \char"0048\cgj\char"0331} & {\itshape \char"0049 \char"0049\cgj\char"0307 \char"0049\cgj\char"0331} & {\itshape \char"004A \char"004A\cgj\char"0307 \char"004A\cgj\char"0331} & {\itshape \char"004B \char"004B\cgj\char"0307 \char"004B\cgj\char"0331} & {\itshape \char"004C \char"004C\cgj\char"0307 \char"004C\cgj\char"0331} & {\itshape \char"004D \char"004D\cgj\char"0307 \char"004D\cgj\char"0331} \\
patch\rule[-0.5em]{0pt}{1.9em} & {\LibertinusSerifPatch\itshape \char"0041 \char"0041\cgj\char"0307 \char"0041\cgj\char"0331} & {\LibertinusSerifPatch\itshape \char"0042 \char"0042\cgj\char"0307 \char"0042\cgj\char"0331} & {\LibertinusSerifPatch\itshape \char"0043 \char"0043\cgj\char"0307 \char"0043\cgj\char"0331} & {\LibertinusSerifPatch\itshape \char"0044 \char"0044\cgj\char"0307 \char"0044\cgj\char"0331} & {\LibertinusSerifPatch\itshape \char"0045 \char"0045\cgj\char"0307 \char"0045\cgj\char"0331} & {\LibertinusSerifPatch\itshape \char"0046 \char"0046\cgj\char"0307 \char"0046\cgj\char"0331} & {\LibertinusSerifPatch\itshape \char"0047 \char"0047\cgj\char"0307 \char"0047\cgj\char"0331} & {\LibertinusSerifPatch\itshape \char"0048 \char"0048\cgj\char"0307 \char"0048\cgj\char"0331} & {\LibertinusSerifPatch\itshape \char"0049 \char"0049\cgj\char"0307 \char"0049\cgj\char"0331} & {\LibertinusSerifPatch\itshape \char"004A \char"004A\cgj\char"0307 \char"004A\cgj\char"0331} & {\LibertinusSerifPatch\itshape \char"004B \char"004B\cgj\char"0307 \char"004B\cgj\char"0331} & {\LibertinusSerifPatch\itshape \char"004C \char"004C\cgj\char"0307 \char"004C\cgj\char"0331} & {\LibertinusSerifPatch\itshape \char"004D \char"004D\cgj\char"0307 \char"004D\cgj\char"0331} \\
a & 505, 850 & 447, 850 & 500, 850 & 396, 850 & 413, 850 & 406, 850 & 513, 850 & 481, 850 & 316, 850 & 392, 850 & 466, 850 & 389, 850 & 557, 850 \\
a? adv & 514, 850 & 459, 850 & 489, 850 & 514, 850 & 444, 850 & 409, 850 & 513, 850 & 517, 850 & 321, 850 & 388, 850 & 499, 850 & 440, 850 & 583, 850 \\
a$\delta$ adv & -9, 0 & -12, 0 & \textbf{11, 0} & -118, 0 & -31, 0 & \textbf{-3, 0} & \textbf{0, 0} & -36, 0 & -5, 0 & \textbf{4, 0} & -33, 0 & -51, 0 & -26, 0 \\
a? ink & 514, 850 & 441, 850 & 486, 850 & 494, 850 & 429, 850 & 400, 850 & 514, 850 & 510, 850 & 315, 850 & 374, 850 & 485, 850 & 422, 850 & 582, 850 \\
a$\delta$ ink & -9, 0 & \textbf{6, 0} & 14, 0 & \textbf{-98, 0} & \textbf{-16, 0} & 6, 0 & -1, 0 & \textbf{-29, 0} & \textbf{1, 0} & 18, 0 & \textbf{-19, 0} & \textbf{-33, 0} & \textbf{-25, 0} \\
b & 308, -108 & 238, -108 & 333, -110 & 279, -107 & 269, -107 & 131, -111 & 345, -110 & 320, -107 & 127, -109 & 109, -319 & 292, -110 & 252, -110 & 363, -110 \\
b? adv & 310, -110 & 255, -110 & 285, -110 & 310, -110 & 240, -110 & 205, -110 & 309, -110 & 313, -110 & 117, -110 & 146, -285 & 295, -110 & 236, -110 & 379, -110 \\
b$\delta$ adv & -2, 2 & -17, 2 & \textbf{48, 0} & -31, 3 & \textbf{29, 3} & -74, -1 & 36, 0 & \textbf{7, 3} & \textbf{10, 1} & -37, -34 & \textbf{-3, 0} & \textbf{16, 0} & -16, 0 \\
b? ink & 310, -110 & 237, -110 & 282, -110 & 290, -110 & 225, -110 & 196, -110 & 310, -110 & 306, -110 & 111, -110 & 132, -285 & 281, -110 & 218, -110 & 378, -110 \\
b$\delta$ ink & -2, 2 & \textbf{1, 2} & 51, 0 & \textbf{-11, 3} & 44, 3 & \textbf{-65, -1} & \textbf{35, 0} & 14, 3 & 16, 1 & \textbf{-23, -34} & 11, 0 & 34, 0 & \textbf{-15, 0} \\
\hline si\rule[-0.5em]{0pt}{1.9em} & {\bfseries\itshape \char"0041 \char"0041\cgj\char"0307 \char"0041\cgj\char"0331} & {\bfseries\itshape \char"0042 \char"0042\cgj\char"0307 \char"0042\cgj\char"0331} & {\bfseries\itshape \char"0043 \char"0043\cgj\char"0307 \char"0043\cgj\char"0331} & {\bfseries\itshape \char"0044 \char"0044\cgj\char"0307 \char"0044\cgj\char"0331} & {\bfseries\itshape \char"0045 \char"0045\cgj\char"0307 \char"0045\cgj\char"0331} & {\bfseries\itshape \char"0046 \char"0046\cgj\char"0307 \char"0046\cgj\char"0331} & {\bfseries\itshape \char"0047 \char"0047\cgj\char"0307 \char"0047\cgj\char"0331} & {\bfseries\itshape \char"0048 \char"0048\cgj\char"0307 \char"0048\cgj\char"0331} & {\bfseries\itshape \char"0049 \char"0049\cgj\char"0307 \char"0049\cgj\char"0331} & {\bfseries\itshape \char"004A \char"004A\cgj\char"0307 \char"004A\cgj\char"0331} & {\bfseries\itshape \char"004B \char"004B\cgj\char"0307 \char"004B\cgj\char"0331} & {\bfseries\itshape \char"004C \char"004C\cgj\char"0307 \char"004C\cgj\char"0331} & {\bfseries\itshape \char"004D \char"004D\cgj\char"0307 \char"004D\cgj\char"0331} \\
patch\rule[-0.5em]{0pt}{1.9em} & {\LibertinusSerifPatch\bfseries\itshape \char"0041 \char"0041\cgj\char"0307 \char"0041\cgj\char"0331} & {\LibertinusSerifPatch\bfseries\itshape \char"0042 \char"0042\cgj\char"0307 \char"0042\cgj\char"0331} & {\LibertinusSerifPatch\bfseries\itshape \char"0043 \char"0043\cgj\char"0307 \char"0043\cgj\char"0331} & {\LibertinusSerifPatch\bfseries\itshape \char"0044 \char"0044\cgj\char"0307 \char"0044\cgj\char"0331} & {\LibertinusSerifPatch\bfseries\itshape \char"0045 \char"0045\cgj\char"0307 \char"0045\cgj\char"0331} & {\LibertinusSerifPatch\bfseries\itshape \char"0046 \char"0046\cgj\char"0307 \char"0046\cgj\char"0331} & {\LibertinusSerifPatch\bfseries\itshape \char"0047 \char"0047\cgj\char"0307 \char"0047\cgj\char"0331} & {\LibertinusSerifPatch\bfseries\itshape \char"0048 \char"0048\cgj\char"0307 \char"0048\cgj\char"0331} & {\LibertinusSerifPatch\bfseries\itshape \char"0049 \char"0049\cgj\char"0307 \char"0049\cgj\char"0331} & {\LibertinusSerifPatch\bfseries\itshape \char"004A \char"004A\cgj\char"0307 \char"004A\cgj\char"0331} & {\LibertinusSerifPatch\bfseries\itshape \char"004B \char"004B\cgj\char"0307 \char"004B\cgj\char"0331} & {\LibertinusSerifPatch\bfseries\itshape \char"004C \char"004C\cgj\char"0307 \char"004C\cgj\char"0331} & {\LibertinusSerifPatch\bfseries\itshape \char"004D \char"004D\cgj\char"0307 \char"004D\cgj\char"0331} \\
a & 354, 850 & 267, 850 & 340, 850 & 306, 850 & 273, 850 & 276, 850 & 365, 850 & 361, 850 & 150, 850 & 432, 850 & 336, 850 & 149, 850 & 431, 850 \\
a? adv & 382, 850 & 319, 850 & 341, 850 & 367, 850 & 292, 850 & 269, 850 & 371, 850 & 396, 850 & 194, 850 & 249, 850 & 369, 850 & 310, 850 & 441, 850 \\
a$\delta$ adv & \textbf{-28, 0} & -52, 0 & \textbf{-1, 0} & -61, 0 & -19, 0 & \textbf{7, 0} & -6, 0 & \textbf{-35, 0} & \textbf{-44, 0} & \textbf{183, 0} & \textbf{-33, 0} & \textbf{-161, 0} & -10, 0 \\
a? ink & 384, 850 & 290, 850 & 336, 850 & 347, 850 & 283, 850 & 265, 850 & 368, 850 & 403, 850 & 202, 850 & 239, 850 & 380, 850 & 314, 850 & 439, 850 \\
a$\delta$ ink & -30, 0 & \textbf{-23, 0} & 4, 0 & \textbf{-41, 0} & \textbf{-10, 0} & 11, 0 & \textbf{-3, 0} & -42, 0 & -52, 0 & 193, 0 & -44, 0 & -165, 0 & \textbf{-8, 0} \\
b & 308, -110 & 268, -110 & 363, -110 & 279, -110 & 279, -110 & 201, -110 & 367, -110 & 365, -110 & 147, -110 & 109, -319 & 342, -110 & 260, -110 & 393, -110 \\
b? adv & 349, -110 & 286, -110 & 308, -110 & 334, -110 & 259, -110 & 236, -110 & 338, -110 & 363, -110 & 161, -110 & 187, -249 & 336, -110 & 277, -110 & 408, -110 \\
b$\delta$ adv & \textbf{-41, 0} & -18, 0 & \textbf{55, 0} & -55, 0 & \textbf{20, 0} & -35, 0 & \textbf{29, 0} & \textbf{2, 0} & \textbf{-14, 0} & -78, -70 & 6, 0 & \textbf{-17, 0} & -15, 0 \\
b? ink & 351, -110 & 257, -110 & 303, -110 & 314, -110 & 250, -110 & 232, -110 & 335, -110 & 370, -110 & 169, -110 & 177, -249 & 347, -110 & 281, -110 & 406, -110 \\
b$\delta$ ink & -43, 0 & \textbf{11, 0} & 60, 0 & \textbf{-35, 0} & 29, 0 & \textbf{-31, 0} & 32, 0 & -5, 0 & -22, 0 & \textbf{-68, -70} & \textbf{-5, 0} & -21, 0 & \textbf{-13, 0} \\

\end{tabular}
}
\end{table}
\clearpage


\begin{table}[p]
\captionsetup{justification=raggedright, singlelinecheck=false}
\caption{Lane comparison for A--Z 2/2}
\label{table:lanes_AZ_2}
\setlength{\tabcolsep}{3.5pt}

{\footnotesize
\begin{tabular}{lrrrrrrrrrrrrr}
hex & 004E & 004F & 0050 & 0051 & 0052 & 0053 & 0054 & 0055 & 0056 & 0057 & 0058 & 0059 & 005A \\
cat & Lu & Lu & Lu & Lu & Lu & Lu & Lu & Lu & Lu & Lu & Lu & Lu & Lu \\
\hline it\rule[-0.5em]{0pt}{1.9em} & {\itshape \char"004E \char"004E\cgj\char"0307 \char"004E\cgj\char"0331} & {\itshape \char"004F \char"004F\cgj\char"0307 \char"004F\cgj\char"0331} & {\itshape \char"0050 \char"0050\cgj\char"0307 \char"0050\cgj\char"0331} & {\itshape \char"0051 \char"0051\cgj\char"0307 \char"0051\cgj\char"0331} & {\itshape \char"0052 \char"0052\cgj\char"0307 \char"0052\cgj\char"0331} & {\itshape \char"0053 \char"0053\cgj\char"0307 \char"0053\cgj\char"0331} & {\itshape \char"0054 \char"0054\cgj\char"0307 \char"0054\cgj\char"0331} & {\itshape \char"0055 \char"0055\cgj\char"0307 \char"0055\cgj\char"0331} & {\itshape \char"0056 \char"0056\cgj\char"0307 \char"0056\cgj\char"0331} & {\itshape \char"0057 \char"0057\cgj\char"0307 \char"0057\cgj\char"0331} & {\itshape \char"0058 \char"0058\cgj\char"0307 \char"0058\cgj\char"0331} & {\itshape \char"0059 \char"0059\cgj\char"0307 \char"0059\cgj\char"0331} & {\itshape \char"005A \char"005A\cgj\char"0307 \char"005A\cgj\char"0331} \\
patch\rule[-0.5em]{0pt}{1.9em} & {\LibertinusSerifPatch\itshape \char"004E \char"004E\cgj\char"0307 \char"004E\cgj\char"0331} & {\LibertinusSerifPatch\itshape \char"004F \char"004F\cgj\char"0307 \char"004F\cgj\char"0331} & {\LibertinusSerifPatch\itshape \char"0050 \char"0050\cgj\char"0307 \char"0050\cgj\char"0331} & {\LibertinusSerifPatch\itshape \char"0051 \char"0051\cgj\char"0307 \char"0051\cgj\char"0331} & {\LibertinusSerifPatch\itshape \char"0052 \char"0052\cgj\char"0307 \char"0052\cgj\char"0331} & {\LibertinusSerifPatch\itshape \char"0053 \char"0053\cgj\char"0307 \char"0053\cgj\char"0331} & {\LibertinusSerifPatch\itshape \char"0054 \char"0054\cgj\char"0307 \char"0054\cgj\char"0331} & {\LibertinusSerifPatch\itshape \char"0055 \char"0055\cgj\char"0307 \char"0055\cgj\char"0331} & {\LibertinusSerifPatch\itshape \char"0056 \char"0056\cgj\char"0307 \char"0056\cgj\char"0331} & {\LibertinusSerifPatch\itshape \char"0057 \char"0057\cgj\char"0307 \char"0057\cgj\char"0331} & {\LibertinusSerifPatch\itshape \char"0058 \char"0058\cgj\char"0307 \char"0058\cgj\char"0331} & {\LibertinusSerifPatch\itshape \char"0059 \char"0059\cgj\char"0307 \char"0059\cgj\char"0331} & {\LibertinusSerifPatch\itshape \char"005A \char"005A\cgj\char"0307 \char"005A\cgj\char"0331} \\
a & 476, 850 & 515, 850 & 426, 850 & 498, 850 & 419, 850 & 402, 850 & 457, 850 & 486, 850 & 481, 850 & 610, 850 & 482, 850 & 435, 850 & 450, 850 \\
a? adv & 514, 850 & 515, 850 & 430, 850 & 515, 850 & 458, 850 & 408, 850 & 453, 850 & 498, 850 & 479, 850 & 635, 850 & 495, 850 & 457, 850 & 470, 850 \\
a$\delta$ adv & -38, 0 & \textbf{0, 0} & \textbf{-4, 0} & \textbf{-17, 0} & \textbf{-39, 0} & -6, 0 & 4, 0 & -12, 0 & \textbf{2, 0} & -25, 0 & -13, 0 & -22, 0 & \textbf{-20, 0} \\
a? ink & 510, 850 & 510, 850 & 417, 850 & 524, 850 & 465, 850 & 403, 850 & 453, 850 & 486, 850 & 464, 850 & 634, 850 & 495, 850 & 438, 850 & 475, 850 \\
a$\delta$ ink & \textbf{-34, 0} & 5, 0 & 9, 0 & -26, 0 & -46, 0 & \textbf{-1, 0} & 4, 0 & \textbf{0, 0} & 17, 0 & \textbf{-24, 0} & -13, 0 & \textbf{-3, 0} & -25, 0 \\
b & 328, -109 & 310, -113 & 140, -110 &  & 248, -111 & 204, -111 & 277, -109 & 308, -109 & 263, -110 & 439, -107 & 287, -107 & 270, -109 & 267, -107 \\
b? adv & 310, -110 & 311, -110 & 226, -110 & 272, -294 & 254, -110 & 204, -110 & 249, -110 & 294, -110 & 275, -110 & 431, -110 & 291, -110 & 253, -110 & 266, -110 \\
b$\delta$ adv & \textbf{18, 1} & \textbf{-1, -3} & -86, 0 &  & \textbf{-6, -1} & \textbf{0, -1} & 28, 1 & \textbf{14, 1} & -12, 0 & \textbf{8, 3} & -4, 3 & \textbf{17, 1} & \textbf{1, 3} \\
b? ink & 306, -110 & 306, -110 & 213, -110 & 281, -294 & 261, -110 & 199, -110 & 249, -110 & 282, -110 & 260, -110 & 430, -110 & 291, -110 & 234, -110 & 271, -110 \\
b$\delta$ ink & 22, 1 & 4, -3 & \textbf{-73, 0} &  & -13, -1 & 5, -1 & 28, 1 & 26, 1 & \textbf{3, 0} & 9, 3 & -4, 3 & 36, 1 & -4, 3 \\
\hline si\rule[-0.5em]{0pt}{1.9em} & {\bfseries\itshape \char"004E \char"004E\cgj\char"0307 \char"004E\cgj\char"0331} & {\bfseries\itshape \char"004F \char"004F\cgj\char"0307 \char"004F\cgj\char"0331} & {\bfseries\itshape \char"0050 \char"0050\cgj\char"0307 \char"0050\cgj\char"0331} & {\bfseries\itshape \char"0051 \char"0051\cgj\char"0307 \char"0051\cgj\char"0331} & {\bfseries\itshape \char"0052 \char"0052\cgj\char"0307 \char"0052\cgj\char"0331} & {\bfseries\itshape \char"0053 \char"0053\cgj\char"0307 \char"0053\cgj\char"0331} & {\bfseries\itshape \char"0054 \char"0054\cgj\char"0307 \char"0054\cgj\char"0331} & {\bfseries\itshape \char"0055 \char"0055\cgj\char"0307 \char"0055\cgj\char"0331} & {\bfseries\itshape \char"0056 \char"0056\cgj\char"0307 \char"0056\cgj\char"0331} & {\bfseries\itshape \char"0057 \char"0057\cgj\char"0307 \char"0057\cgj\char"0331} & {\bfseries\itshape \char"0058 \char"0058\cgj\char"0307 \char"0058\cgj\char"0331} & {\bfseries\itshape \char"0059 \char"0059\cgj\char"0307 \char"0059\cgj\char"0331} & {\bfseries\itshape \char"005A \char"005A\cgj\char"0307 \char"005A\cgj\char"0331} \\
patch\rule[-0.5em]{0pt}{1.9em} & {\LibertinusSerifPatch\bfseries\itshape \char"004E \char"004E\cgj\char"0307 \char"004E\cgj\char"0331} & {\LibertinusSerifPatch\bfseries\itshape \char"004F \char"004F\cgj\char"0307 \char"004F\cgj\char"0331} & {\LibertinusSerifPatch\bfseries\itshape \char"0050 \char"0050\cgj\char"0307 \char"0050\cgj\char"0331} & {\LibertinusSerifPatch\bfseries\itshape \char"0051 \char"0051\cgj\char"0307 \char"0051\cgj\char"0331} & {\LibertinusSerifPatch\bfseries\itshape \char"0052 \char"0052\cgj\char"0307 \char"0052\cgj\char"0331} & {\LibertinusSerifPatch\bfseries\itshape \char"0053 \char"0053\cgj\char"0307 \char"0053\cgj\char"0331} & {\LibertinusSerifPatch\bfseries\itshape \char"0054 \char"0054\cgj\char"0307 \char"0054\cgj\char"0331} & {\LibertinusSerifPatch\bfseries\itshape \char"0055 \char"0055\cgj\char"0307 \char"0055\cgj\char"0331} & {\LibertinusSerifPatch\bfseries\itshape \char"0056 \char"0056\cgj\char"0307 \char"0056\cgj\char"0331} & {\LibertinusSerifPatch\bfseries\itshape \char"0057 \char"0057\cgj\char"0307 \char"0057\cgj\char"0331} & {\LibertinusSerifPatch\bfseries\itshape \char"0058 \char"0058\cgj\char"0307 \char"0058\cgj\char"0331} & {\LibertinusSerifPatch\bfseries\itshape \char"0059 \char"0059\cgj\char"0307 \char"0059\cgj\char"0331} & {\LibertinusSerifPatch\bfseries\itshape \char"005A \char"005A\cgj\char"0307 \char"005A\cgj\char"0331} \\
a & 346, 850 & 335, 850 & 316, 850 & 348, 850 & 269, 850 & 232, 850 & 307, 850 & 346, 850 & 336, 850 & 450, 850 & 332, 850 & 295, 850 & 320, 850 \\
a? adv & 394, 850 & 390, 850 & 284, 850 & 390, 850 & 315, 850 & 250, 850 & 309, 850 & 364, 850 & 328, 850 & 501, 850 & 362, 850 & 319, 850 & 322, 850 \\
a$\delta$ adv & -48, 0 & \textbf{-55, 0} & \textbf{32, 0} & \textbf{-42, 0} & \textbf{-46, 0} & -18, 0 & \textbf{-2, 0} & \textbf{-18, 0} & \textbf{8, 0} & -51, 0 & \textbf{-30, 0} & -24, 0 & \textbf{-2, 0} \\
a? ink & 380, 850 & 393, 850 & 271, 850 & 393, 850 & 322, 850 & 245, 850 & 296, 850 & 365, 850 & 312, 850 & 475, 850 & 363, 850 & 302, 850 & 317, 850 \\
a$\delta$ ink & \textbf{-34, 0} & -58, 0 & 45, 0 & -45, 0 & -53, 0 & \textbf{-13, 0} & 11, 0 & -19, 0 & 24, 0 & \textbf{-25, 0} & -31, 0 & \textbf{-7, 0} & 3, 0 \\
b & 348, -110 & 350, -110 & 240, -110 &  & 258, -110 & 234, -110 & 300, -110 & 328, -110 & 318, -110 & 479, -110 & 307, -110 & 295, -110 & 287, -110 \\
b? adv & 361, -110 & 357, -110 & 251, -110 & 314, -319 & 282, -110 & 217, -110 & 276, -110 & 331, -110 & 295, -110 & 468, -110 & 329, -110 & 286, -110 & 289, -110 \\
b$\delta$ adv & -13, 0 & \textbf{-7, 0} & -11, 0 &  & \textbf{-24, 0} & \textbf{17, 0} & \textbf{24, 0} & \textbf{-3, 0} & \textbf{23, 0} & \textbf{11, 0} & \textbf{-22, 0} & \textbf{9, 0} & \textbf{-2, 0} \\
b? ink & 347, -110 & 360, -110 & 238, -110 & 317, -319 & 289, -110 & 212, -110 & 263, -110 & 332, -110 & 279, -110 & 442, -110 & 330, -110 & 269, -110 & 284, -110 \\
b$\delta$ ink & \textbf{1, 0} & -10, 0 & \textbf{2, 0} &  & -31, 0 & 22, 0 & 37, 0 & -4, 0 & 39, 0 & 37, 0 & -23, 0 & 26, 0 & 3, 0 \\

\end{tabular}
}
\end{table}
\clearpage



\begin{table}[p]
\captionsetup{justification=raggedright, singlelinecheck=false}
\caption{Lane comparison for a--z 1/2}
\label{table:lanes_az_1}
\setlength{\tabcolsep}{3.5pt}

{\footnotesize
\begin{tabular}{lrrrrrrrrrrrrr}
hex & 0061 & 0062 & 0063 & 0064 & 0065 & 0066 & 0067 & 0068 & 0069 & 006A & 006B & 006C & 006D \\
cat & Ll & Ll & Ll & Ll & Ll & Ll & Ll & Ll & Ll & Ll & Ll & Ll & Ll \\
\hline it\rule[-0.5em]{0pt}{1.9em} & {\itshape \char"0061 \char"0061\cgj\char"0307 \char"0061\cgj\char"0331} & {\itshape \char"0062 \char"0062\cgj\char"0307 \char"0062\cgj\char"0331} & {\itshape \char"0063 \char"0063\cgj\char"0307 \char"0063\cgj\char"0331} & {\itshape \char"0064 \char"0064\cgj\char"0307 \char"0064\cgj\char"0331} & {\itshape \char"0065 \char"0065\cgj\char"0307 \char"0065\cgj\char"0331} & {\itshape \char"0066 \char"0066\cgj\char"0307 \char"0066\cgj\char"0331} & {\itshape \char"0067 \char"0067\cgj\char"0307 \char"0067\cgj\char"0331} & {\itshape \char"0068 \char"0068\cgj\char"0307 \char"0068\cgj\char"0331} & {\itshape \char"0069 \char"0069\cgj\char"0307 \char"0069\cgj\char"0331} & {\itshape \char"006A \char"006A\cgj\char"0307 \char"006A\cgj\char"0331} & {\itshape \char"006B \char"006B\cgj\char"0307 \char"006B\cgj\char"0331} & {\itshape \char"006C \char"006C\cgj\char"0307 \char"006C\cgj\char"0331} & {\itshape \char"006D \char"006D\cgj\char"0307 \char"006D\cgj\char"0331} \\
patch\rule[-0.5em]{0pt}{1.9em} & {\LibertinusSerifPatch\itshape \char"0061 \char"0061\cgj\char"0307 \char"0061\cgj\char"0331} & {\LibertinusSerifPatch\itshape \char"0062 \char"0062\cgj\char"0307 \char"0062\cgj\char"0331} & {\LibertinusSerifPatch\itshape \char"0063 \char"0063\cgj\char"0307 \char"0063\cgj\char"0331} & {\LibertinusSerifPatch\itshape \char"0064 \char"0064\cgj\char"0307 \char"0064\cgj\char"0331} & {\LibertinusSerifPatch\itshape \char"0065 \char"0065\cgj\char"0307 \char"0065\cgj\char"0331} & {\LibertinusSerifPatch\itshape \char"0066 \char"0066\cgj\char"0307 \char"0066\cgj\char"0331} & {\LibertinusSerifPatch\itshape \char"0067 \char"0067\cgj\char"0307 \char"0067\cgj\char"0331} & {\LibertinusSerifPatch\itshape \char"0068 \char"0068\cgj\char"0307 \char"0068\cgj\char"0331} & {\LibertinusSerifPatch\itshape \char"0069 \char"0069\cgj\char"0307 \char"0069\cgj\char"0331} & {\LibertinusSerifPatch\itshape \char"006A \char"006A\cgj\char"0307 \char"006A\cgj\char"0331} & {\LibertinusSerifPatch\itshape \char"006B \char"006B\cgj\char"0307 \char"006B\cgj\char"0331} & {\LibertinusSerifPatch\itshape \char"006C \char"006C\cgj\char"0307 \char"006C\cgj\char"0331} & {\LibertinusSerifPatch\itshape \char"006D \char"006D\cgj\char"0307 \char"006D\cgj\char"0331} \\
a & 341, 647 &  & 343, 645 & 218, 645 & 332, 645 &  & 360, 645 & 409, 645 &  &  & 405, 646 & 292, 890 & 522, 645 \\
a? adv & 380, 645 & 428, 890 & 332, 645 & 434, 890 & 338, 645 & 346, 890 & 376, 645 & 449, 890 & 310, 810 & 302, 810 & 432, 890 & 322, 890 & 529, 645 \\
a$\delta$ adv & \textbf{-39, 2} &  & 11, 0 & \textbf{-216, -245} & \textbf{-6, 0} &  & \textbf{-16, 0} & \textbf{-40, -245} &  &  & -27, -244 & \textbf{-30, 0} & \textbf{-7, 0} \\
a? ink & 396, 645 & 411, 890 & 344, 645 & 451, 890 & 345, 645 & 365, 890 & 380, 645 & 455, 890 & 321, 810 & 240, 810 & 430, 890 & 327, 890 & 539, 645 \\
a$\delta$ ink & -55, 2 &  & \textbf{-1, 0} & -233, -245 & -13, 0 &  & -20, 0 & -46, -245 &  &  & \textbf{-25, -244} & -35, 0 & -17, 0 \\
b & 174, -108 & 214, -108 & 163, -113 & 221, -107 & 158, -108 & 222, -111 & 197, -319 & 229, -111 & 136, -107 & 0, -319 & 223, -107 & 115, -107 & 382, -105 \\
b? adv & 220, -110 & 216, -110 & 171, -110 & 221, -110 & 177, -110 & 89, -319 & 171, -319 & 236, -110 & 115, -110 & 70, -281 & 220, -110 & 110, -110 & 368, -110 \\
b$\delta$ adv & \textbf{-46, 2} & \textbf{-2, 2} & \textbf{-8, -3} & \textbf{0, 3} & \textbf{-19, 2} & 133, 208 & 26, 0 & \textbf{-7, -1} & 21, 3 & -70, -38 & \textbf{3, 3} & 5, 3 & 14, 5 \\
b? ink & 236, -110 & 199, -110 & 184, -110 & 239, -110 & 185, -110 & 108, -319 & 175, -319 & 243, -110 & 126, -110 & 8, -281 & 218, -110 & 115, -110 & 379, -110 \\
b$\delta$ ink & -62, 2 & 15, 2 & -21, -3 & -18, 3 & -27, 2 & \textbf{114, 208} & \textbf{22, 0} & -14, -1 & \textbf{10, 3} & \textbf{-8, -38} & 5, 3 & \textbf{0, 3} & \textbf{3, 5} \\
\hline si\rule[-0.5em]{0pt}{1.9em} & {\bfseries\itshape \char"0061 \char"0061\cgj\char"0307 \char"0061\cgj\char"0331} & {\bfseries\itshape \char"0062 \char"0062\cgj\char"0307 \char"0062\cgj\char"0331} & {\bfseries\itshape \char"0063 \char"0063\cgj\char"0307 \char"0063\cgj\char"0331} & {\bfseries\itshape \char"0064 \char"0064\cgj\char"0307 \char"0064\cgj\char"0331} & {\bfseries\itshape \char"0065 \char"0065\cgj\char"0307 \char"0065\cgj\char"0331} & {\bfseries\itshape \char"0066 \char"0066\cgj\char"0307 \char"0066\cgj\char"0331} & {\bfseries\itshape \char"0067 \char"0067\cgj\char"0307 \char"0067\cgj\char"0331} & {\bfseries\itshape \char"0068 \char"0068\cgj\char"0307 \char"0068\cgj\char"0331} & {\bfseries\itshape \char"0069 \char"0069\cgj\char"0307 \char"0069\cgj\char"0331} & {\bfseries\itshape \char"006A \char"006A\cgj\char"0307 \char"006A\cgj\char"0331} & {\bfseries\itshape \char"006B \char"006B\cgj\char"0307 \char"006B\cgj\char"0331} & {\bfseries\itshape \char"006C \char"006C\cgj\char"0307 \char"006C\cgj\char"0331} & {\bfseries\itshape \char"006D \char"006D\cgj\char"0307 \char"006D\cgj\char"0331} \\
patch\rule[-0.5em]{0pt}{1.9em} & {\LibertinusSerifPatch\bfseries\itshape \char"0061 \char"0061\cgj\char"0307 \char"0061\cgj\char"0331} & {\LibertinusSerifPatch\bfseries\itshape \char"0062 \char"0062\cgj\char"0307 \char"0062\cgj\char"0331} & {\LibertinusSerifPatch\bfseries\itshape \char"0063 \char"0063\cgj\char"0307 \char"0063\cgj\char"0331} & {\LibertinusSerifPatch\bfseries\itshape \char"0064 \char"0064\cgj\char"0307 \char"0064\cgj\char"0331} & {\LibertinusSerifPatch\bfseries\itshape \char"0065 \char"0065\cgj\char"0307 \char"0065\cgj\char"0331} & {\LibertinusSerifPatch\bfseries\itshape \char"0066 \char"0066\cgj\char"0307 \char"0066\cgj\char"0331} & {\LibertinusSerifPatch\bfseries\itshape \char"0067 \char"0067\cgj\char"0307 \char"0067\cgj\char"0331} & {\LibertinusSerifPatch\bfseries\itshape \char"0068 \char"0068\cgj\char"0307 \char"0068\cgj\char"0331} & {\LibertinusSerifPatch\bfseries\itshape \char"0069 \char"0069\cgj\char"0307 \char"0069\cgj\char"0331} & {\LibertinusSerifPatch\bfseries\itshape \char"006A \char"006A\cgj\char"0307 \char"006A\cgj\char"0331} & {\LibertinusSerifPatch\bfseries\itshape \char"006B \char"006B\cgj\char"0307 \char"006B\cgj\char"0331} & {\LibertinusSerifPatch\bfseries\itshape \char"006C \char"006C\cgj\char"0307 \char"006C\cgj\char"0331} & {\LibertinusSerifPatch\bfseries\itshape \char"006D \char"006D\cgj\char"0307 \char"006D\cgj\char"0331} \\
a & 379, 645 &  & 242, 645 & 218, 645 & 243, 645 &  & 291, 645 & 469, 645 & 137, 795 &  & 306, 784 & 174, 890 & 392, 645 \\
a? adv & 251, 645 & 290, 890 & 203, 645 & 309, 890 & 218, 645 & 216, 890 & 220, 645 & 331, 890 & 180, 823 & 192, 823 & 316, 890 & 194, 890 & 407, 645 \\
a$\delta$ adv & 128, 0 &  & 39, 0 & \textbf{-91, -245} & \textbf{25, 0} &  & 71, 0 & 138, -245 & \textbf{-43, -28} &  & -10, -106 & -20, 0 & \textbf{-15, 0} \\
a? ink & 274, 645 & 281, 890 & 206, 645 & 317, 890 & 217, 645 & 206, 890 & 233, 645 & 331, 890 & 191, 823 & 113, 823 & 314, 890 & 193, 890 & 412, 645 \\
a$\delta$ ink & \textbf{105, 0} &  & \textbf{36, 0} & -99, -245 & 26, 0 &  & \textbf{58, 0} & 138, -245 & -54, -28 &  & \textbf{-8, -106} & \textbf{-19, 0} & -20, 0 \\
b & 196, -110 & 234, -110 & 236, -110 & 221, -110 & 234, -110 & 132, -110 & 199, -319 & 272, -110 & 136, -110 & 67, -319 & 290, -214 & 135, -110 & 400, -110 \\
b? adv & 260, -110 & 248, -110 & 212, -110 & 268, -110 & 227, -110 & 132, -319 & 186, -319 & 290, -110 & 153, -110 & 127, -292 & 275, -110 & 153, -110 & 415, -110 \\
b$\delta$ adv & \textbf{-64, 0} & -14, 0 & 24, 0 & \textbf{-47, 0} & \textbf{7, 0} & \textbf{0, 209} & 13, 0 & -18, 0 & \textbf{-17, 0} & -60, -27 & \textbf{15, -104} & -18, 0 & \textbf{-15, 0} \\
b? ink & 283, -110 & 240, -110 & 215, -110 & 276, -110 & 226, -110 & 122, -319 & 199, -319 & 290, -110 & 164, -110 & 49, -292 & 273, -110 & 152, -110 & 421, -110 \\
b$\delta$ ink & -87, 0 & \textbf{-6, 0} & \textbf{21, 0} & -55, 0 & 8, 0 & 10, 209 & \textbf{0, 0} & -18, 0 & -28, 0 & \textbf{18, -27} & 17, -104 & \textbf{-17, 0} & -21, 0 \\

\end{tabular}
}
\end{table}
\clearpage


\begin{table}[p]
\captionsetup{justification=raggedright, singlelinecheck=false}
\caption{Lane comparison for a--z 2/2}
\label{table:lanes_az_2}
\setlength{\tabcolsep}{3.5pt}

{\footnotesize
\begin{tabular}{lrrrrrrrrrrrrr}
hex & 006E & 006F & 0070 & 0071 & 0072 & 0073 & 0074 & 0075 & 0076 & 0077 & 0078 & 0079 & 007A \\
cat & Ll & Ll & Ll & Ll & Ll & Ll & Ll & Ll & Ll & Ll & Ll & Ll & Ll \\
\hline it\rule[-0.5em]{0pt}{1.9em} & {\itshape \char"006E \char"006E\cgj\char"0307 \char"006E\cgj\char"0331} & {\itshape \char"006F \char"006F\cgj\char"0307 \char"006F\cgj\char"0331} & {\itshape \char"0070 \char"0070\cgj\char"0307 \char"0070\cgj\char"0331} & {\itshape \char"0071 \char"0071\cgj\char"0307 \char"0071\cgj\char"0331} & {\itshape \char"0072 \char"0072\cgj\char"0307 \char"0072\cgj\char"0331} & {\itshape \char"0073 \char"0073\cgj\char"0307 \char"0073\cgj\char"0331} & {\itshape \char"0074 \char"0074\cgj\char"0307 \char"0074\cgj\char"0331} & {\itshape \char"0075 \char"0075\cgj\char"0307 \char"0075\cgj\char"0331} & {\itshape \char"0076 \char"0076\cgj\char"0307 \char"0076\cgj\char"0331} & {\itshape \char"0077 \char"0077\cgj\char"0307 \char"0077\cgj\char"0331} & {\itshape \char"0078 \char"0078\cgj\char"0307 \char"0078\cgj\char"0331} & {\itshape \char"0079 \char"0079\cgj\char"0307 \char"0079\cgj\char"0331} & {\itshape \char"007A \char"007A\cgj\char"0307 \char"007A\cgj\char"0331} \\
patch\rule[-0.5em]{0pt}{1.9em} & {\LibertinusSerifPatch\itshape \char"006E \char"006E\cgj\char"0307 \char"006E\cgj\char"0331} & {\LibertinusSerifPatch\itshape \char"006F \char"006F\cgj\char"0307 \char"006F\cgj\char"0331} & {\LibertinusSerifPatch\itshape \char"0070 \char"0070\cgj\char"0307 \char"0070\cgj\char"0331} & {\LibertinusSerifPatch\itshape \char"0071 \char"0071\cgj\char"0307 \char"0071\cgj\char"0331} & {\LibertinusSerifPatch\itshape \char"0072 \char"0072\cgj\char"0307 \char"0072\cgj\char"0331} & {\LibertinusSerifPatch\itshape \char"0073 \char"0073\cgj\char"0307 \char"0073\cgj\char"0331} & {\LibertinusSerifPatch\itshape \char"0074 \char"0074\cgj\char"0307 \char"0074\cgj\char"0331} & {\LibertinusSerifPatch\itshape \char"0075 \char"0075\cgj\char"0307 \char"0075\cgj\char"0331} & {\LibertinusSerifPatch\itshape \char"0076 \char"0076\cgj\char"0307 \char"0076\cgj\char"0331} & {\LibertinusSerifPatch\itshape \char"0077 \char"0077\cgj\char"0307 \char"0077\cgj\char"0331} & {\LibertinusSerifPatch\itshape \char"0078 \char"0078\cgj\char"0307 \char"0078\cgj\char"0331} & {\LibertinusSerifPatch\itshape \char"0079 \char"0079\cgj\char"0307 \char"0079\cgj\char"0331} & {\LibertinusSerifPatch\itshape \char"007A \char"007A\cgj\char"0307 \char"007A\cgj\char"0331} \\
a & 394, 646 & 359, 645 & 411, 646 & 369, 645 & 312, 646 & 327, 645 &  & 386, 646 &  & 348, 646 & 366, 646 & 367, 644 & 360, 646 \\
a? adv & 396, 645 & 361, 645 & 394, 705 & 383, 645 & 316, 645 & 314, 645 & 318, 775 & 398, 645 & 373, 645 & 481, 645 & 375, 645 & 389, 645 & 355, 645 \\
a$\delta$ adv & -2, 1 & \textbf{-2, 0} & \textbf{17, -59} & \textbf{-14, 0} & \textbf{-4, 1} & 13, 0 &  & \textbf{-12, 1} &  & -133, 1 & -9, 1 & -22, -1 & 5, 1 \\
a? ink & 396, 645 & 362, 645 & 382, 705 & 398, 645 & 317, 645 & 321, 645 & 312, 775 & 402, 645 & 363, 645 & 471, 645 & 371, 645 & 378, 645 & 358, 645 \\
a$\delta$ ink & -2, 1 & -3, 0 & 29, -59 & -29, 0 & -5, 1 & \textbf{6, 0} &  & -16, 1 &  & \textbf{-123, 1} & \textbf{-5, 1} & \textbf{-11, -1} & \textbf{2, 1} \\
b & 235, -111 & 186, -107 & 248, -110 & 169, -110 & 113, -110 & 169, -111 & 133, -107 & 226, -111 &  & 355, -110 & 224, -110 & 368, -107 & 187, -111 \\
b? adv & 236, -110 & 200, -110 & 177, -319 & 178, -319 & 155, -110 & 153, -110 & 130, -110 & 237, -110 & 213, -110 & 321, -110 & 214, -110 & 184, -319 & 195, -110 \\
b$\delta$ adv & -1, -1 & \textbf{-14, 3} & \textbf{71, 209} & \textbf{-9, 209} & \textbf{-42, 0} & 16, -1 & \textbf{3, 3} & \textbf{-11, -1} &  & \textbf{34, 0} & \textbf{10, 0} & \textbf{184, 212} & \textbf{-8, -1} \\
b? ink & 236, -110 & 202, -110 & 164, -319 & 193, -319 & 157, -110 & 161, -110 & 124, -110 & 242, -110 & 203, -110 & 311, -110 & 211, -110 & 173, -319 & 198, -110 \\
b$\delta$ ink & -1, -1 & -16, 3 & 84, 209 & -24, 209 & -44, 0 & \textbf{8, -1} & 9, 3 & -16, -1 &  & 44, 0 & 13, 0 & 195, 212 & -11, -1 \\
\hline si\rule[-0.5em]{0pt}{1.9em} & {\bfseries\itshape \char"006E \char"006E\cgj\char"0307 \char"006E\cgj\char"0331} & {\bfseries\itshape \char"006F \char"006F\cgj\char"0307 \char"006F\cgj\char"0331} & {\bfseries\itshape \char"0070 \char"0070\cgj\char"0307 \char"0070\cgj\char"0331} & {\bfseries\itshape \char"0071 \char"0071\cgj\char"0307 \char"0071\cgj\char"0331} & {\bfseries\itshape \char"0072 \char"0072\cgj\char"0307 \char"0072\cgj\char"0331} & {\bfseries\itshape \char"0073 \char"0073\cgj\char"0307 \char"0073\cgj\char"0331} & {\bfseries\itshape \char"0074 \char"0074\cgj\char"0307 \char"0074\cgj\char"0331} & {\bfseries\itshape \char"0075 \char"0075\cgj\char"0307 \char"0075\cgj\char"0331} & {\bfseries\itshape \char"0076 \char"0076\cgj\char"0307 \char"0076\cgj\char"0331} & {\bfseries\itshape \char"0077 \char"0077\cgj\char"0307 \char"0077\cgj\char"0331} & {\bfseries\itshape \char"0078 \char"0078\cgj\char"0307 \char"0078\cgj\char"0331} & {\bfseries\itshape \char"0079 \char"0079\cgj\char"0307 \char"0079\cgj\char"0331} & {\bfseries\itshape \char"007A \char"007A\cgj\char"0307 \char"007A\cgj\char"0331} \\
patch\rule[-0.5em]{0pt}{1.9em} & {\LibertinusSerifPatch\bfseries\itshape \char"006E \char"006E\cgj\char"0307 \char"006E\cgj\char"0331} & {\LibertinusSerifPatch\bfseries\itshape \char"006F \char"006F\cgj\char"0307 \char"006F\cgj\char"0331} & {\LibertinusSerifPatch\bfseries\itshape \char"0070 \char"0070\cgj\char"0307 \char"0070\cgj\char"0331} & {\LibertinusSerifPatch\bfseries\itshape \char"0071 \char"0071\cgj\char"0307 \char"0071\cgj\char"0331} & {\LibertinusSerifPatch\bfseries\itshape \char"0072 \char"0072\cgj\char"0307 \char"0072\cgj\char"0331} & {\LibertinusSerifPatch\bfseries\itshape \char"0073 \char"0073\cgj\char"0307 \char"0073\cgj\char"0331} & {\LibertinusSerifPatch\bfseries\itshape \char"0074 \char"0074\cgj\char"0307 \char"0074\cgj\char"0331} & {\LibertinusSerifPatch\bfseries\itshape \char"0075 \char"0075\cgj\char"0307 \char"0075\cgj\char"0331} & {\LibertinusSerifPatch\bfseries\itshape \char"0076 \char"0076\cgj\char"0307 \char"0076\cgj\char"0331} & {\LibertinusSerifPatch\bfseries\itshape \char"0077 \char"0077\cgj\char"0307 \char"0077\cgj\char"0331} & {\LibertinusSerifPatch\bfseries\itshape \char"0078 \char"0078\cgj\char"0307 \char"0078\cgj\char"0331} & {\LibertinusSerifPatch\bfseries\itshape \char"0079 \char"0079\cgj\char"0307 \char"0079\cgj\char"0331} & {\LibertinusSerifPatch\bfseries\itshape \char"007A \char"007A\cgj\char"0307 \char"007A\cgj\char"0331} \\
a & 264, 646 & 238, 645 & 271, 646 & 238, 645 & 191, 646 & 326, 645 &  & 255, 645 & 256, 646 & 345, 646 &  & 262, 645 & 220, 646 \\
a? adv & 284, 645 & 235, 645 & 257, 679 & 267, 645 & 203, 645 & 190, 645 & 189, 765 & 262, 645 & 220, 645 & 332, 645 & 240, 645 & 236, 645 & 210, 645 \\
a$\delta$ adv & -20, 1 & \textbf{3, 0} & \textbf{14, -33} & \textbf{-29, 0} & -12, 1 & 136, 0 &  & -7, 0 & \textbf{36, 1} & \textbf{13, 1} &  & \textbf{26, 0} & 10, 1 \\
a? ink & 283, 645 & 231, 645 & 239, 679 & 269, 645 & 202, 645 & 196, 645 & 184, 765 & 261, 645 & 218, 645 & 326, 645 & 238, 645 & 227, 645 & 213, 645 \\
a$\delta$ ink & \textbf{-19, 1} & 7, 0 & 32, -33 & -31, 0 & \textbf{-11, 1} & \textbf{130, 0} &  & \textbf{-6, 0} & 38, 1 & 19, 1 &  & 35, 0 & \textbf{7, 1} \\
b & 265, -110 & 243, -110 & 179, -319 & 348, -319 & 135, -110 & 196, -110 & 133, -107 & 239, -110 & 241, -110 & 352, -110 &  & 242, -319 & 207, -110 \\
b? adv & 293, -110 & 243, -110 & 216, -319 & 233, -319 & 212, -110 & 198, -110 & 173, -110 & 270, -110 & 228, -110 & 341, -110 & 248, -110 & 202, -319 & 218, -110 \\
b$\delta$ adv & -28, 0 & \textbf{0, 0} & -37, 0 & 115, 0 & -77, 0 & \textbf{-2, 0} & -40, 3 & -31, 0 & \textbf{13, 0} & \textbf{11, 0} &  & \textbf{40, 0} & \textbf{-11, 0} \\
b? ink & 292, -110 & 240, -110 & 198, -319 & 235, -319 & 211, -110 & 205, -110 & 168, -110 & 270, -110 & 227, -110 & 335, -110 & 247, -110 & 193, -319 & 222, -110 \\
b$\delta$ ink & \textbf{-27, 0} & 3, 0 & \textbf{-19, 0} & \textbf{113, 0} & \textbf{-76, 0} & -9, 0 & \textbf{-35, 3} & -31, 0 & 14, 0 & 17, 0 &  & 49, 0 & -15, 0 \\

\end{tabular}
}
\end{table}
\clearpage



\begin{table}[p]
\captionsetup{justification=raggedright, singlelinecheck=false}
\caption{Lane comparison for Latin and IPA 1/8}
\label{table:lanes_latinIPA_1}
\setlength{\tabcolsep}{3.5pt}

{\footnotesize
\begin{tabular}{lrrrrrrrrrrrrrrrrr}
hex & 00C6 & 00D0 & 00D8 & 00DE & 00DF & 00E6 & 00F0 & 00F8 & 00FE & 0110 & 0111 & 0126 & 0127 & 0131 & 0141 & 0142 & 014A \\
cat & Lu & Lu & Lu & Lu & Ll & Ll & Ll & Ll & Ll & Lu & Ll & Lu & Ll & Ll & Lu & Ll & Lu \\
\hline it\rule[-0.5em]{0pt}{1.9em} & {\itshape \char"00C6 \char"00C6\cgj\char"0307 \char"00C6\cgj\char"0331} & {\itshape \char"00D0 \char"00D0\cgj\char"0307 \char"00D0\cgj\char"0331} & {\itshape \char"00D8 \char"00D8\cgj\char"0307 \char"00D8\cgj\char"0331} & {\itshape \char"00DE \char"00DE\cgj\char"0307 \char"00DE\cgj\char"0331} & {\itshape \char"00DF \char"00DF\cgj\char"0307 \char"00DF\cgj\char"0331} & {\itshape \char"00E6 \char"00E6\cgj\char"0307 \char"00E6\cgj\char"0331} & {\itshape \char"00F0 \char"00F0\cgj\char"0307 \char"00F0\cgj\char"0331} & {\itshape \char"00F8 \char"00F8\cgj\char"0307 \char"00F8\cgj\char"0331} & {\itshape \char"00FE \char"00FE\cgj\char"0307 \char"00FE\cgj\char"0331} & {\itshape \char"0110 \char"0110\cgj\char"0307 \char"0110\cgj\char"0331} & {\itshape \char"0111 \char"0111\cgj\char"0307 \char"0111\cgj\char"0331} & {\itshape \char"0126 \char"0126\cgj\char"0307 \char"0126\cgj\char"0331} & {\itshape \char"0127 \char"0127\cgj\char"0307 \char"0127\cgj\char"0331} & {\itshape \char"0131 \char"0131\cgj\char"0307 \char"0131\cgj\char"0331} & {\itshape \char"0141 \char"0141\cgj\char"0307 \char"0141\cgj\char"0331} & {\itshape \char"0142 \char"0142\cgj\char"0307 \char"0142\cgj\char"0331} & {\itshape \char"014A \char"014A\cgj\char"0307 \char"014A\cgj\char"0331} \\
patch\rule[-0.5em]{0pt}{1.9em} & {\LibertinusSerifPatch\itshape \char"00C6 \char"00C6\cgj\char"0307 \char"00C6\cgj\char"0331} & {\LibertinusSerifPatch\itshape \char"00D0 \char"00D0\cgj\char"0307 \char"00D0\cgj\char"0331} & {\LibertinusSerifPatch\itshape \char"00D8 \char"00D8\cgj\char"0307 \char"00D8\cgj\char"0331} & {\LibertinusSerifPatch\itshape \char"00DE \char"00DE\cgj\char"0307 \char"00DE\cgj\char"0331} & {\LibertinusSerifPatch\itshape \char"00DF \char"00DF\cgj\char"0307 \char"00DF\cgj\char"0331} & {\LibertinusSerifPatch\itshape \char"00E6 \char"00E6\cgj\char"0307 \char"00E6\cgj\char"0331} & {\LibertinusSerifPatch\itshape \char"00F0 \char"00F0\cgj\char"0307 \char"00F0\cgj\char"0331} & {\LibertinusSerifPatch\itshape \char"00F8 \char"00F8\cgj\char"0307 \char"00F8\cgj\char"0331} & {\LibertinusSerifPatch\itshape \char"00FE \char"00FE\cgj\char"0307 \char"00FE\cgj\char"0331} & {\LibertinusSerifPatch\itshape \char"0110 \char"0110\cgj\char"0307 \char"0110\cgj\char"0331} & {\LibertinusSerifPatch\itshape \char"0111 \char"0111\cgj\char"0307 \char"0111\cgj\char"0331} & {\LibertinusSerifPatch\itshape \char"0126 \char"0126\cgj\char"0307 \char"0126\cgj\char"0331} & {\LibertinusSerifPatch\itshape \char"0127 \char"0127\cgj\char"0307 \char"0127\cgj\char"0331} & {\LibertinusSerifPatch\itshape \char"0131 \char"0131\cgj\char"0307 \char"0131\cgj\char"0331} & {\LibertinusSerifPatch\itshape \char"0141 \char"0141\cgj\char"0307 \char"0141\cgj\char"0331} & {\LibertinusSerifPatch\itshape \char"0142 \char"0142\cgj\char"0307 \char"0142\cgj\char"0331} & {\LibertinusSerifPatch\itshape \char"014A \char"014A\cgj\char"0307 \char"014A\cgj\char"0331} \\
a & 615, 850 &  &  &  &  & 480, 645 &  &  & 278, 646 &  &  &  &  & 267, 645 &  &  & 346, 804 \\
a? adv & 595, 850 & 514, 850 & 515, 850 & 431, 850 & 461, 890 & 481, 645 & 429, 896 & 370, 691 & 439, 890 & 514, 850 & 434, 890 & 517, 850 & 449, 890 & 275, 645 & 440, 850 & 322, 890 & 495, 850 \\
a$\delta$ adv & \textbf{20, 0} &  &  &  &  & \textbf{-1, 0} &  &  & -161, -244 &  &  &  &  & \textbf{-8, 0} &  &  & -149, -46 \\
a? ink & 581, 850 & 494, 850 & 510, 850 & 411, 850 & 418, 890 & 485, 645 & 431, 896 & 372, 691 & 429, 890 & 494, 850 & 451, 890 & 510, 850 & 455, 890 & 286, 645 & 422, 850 & 322, 890 & 471, 850 \\
a$\delta$ ink & 34, 0 &  &  &  &  & -5, 0 &  &  & \textbf{-151, -244} &  &  &  &  & -19, 0 &  &  & \textbf{-125, -46} \\
b & 549, -107 &  & 310, -113 &  &  & 327, -109 &  & 186, -107 &  &  &  &  &  & 139, -108 &  &  & 364, -109 \\
b? adv & 391, -110 & 310, -110 & 302, -151 & 227, -110 & 204, -319 & 320, -110 & 216, -110 & 190, -156 & 181, -324 & 310, -110 & 221, -110 & 313, -110 & 236, -110 & 115, -110 & 236, -110 & 110, -110 & 253, -285 \\
b$\delta$ adv & \textbf{158, 3} &  & \textbf{8, 38} &  &  & 7, 1 &  & \textbf{-4, 49} &  &  &  &  &  & 24, 2 &  &  & \textbf{111, 176} \\
b? ink & 377, -110 & 290, -110 & 297, -151 & 207, -110 & 161, -319 & 325, -110 & 218, -110 & 192, -156 & 171, -324 & 290, -110 & 239, -110 & 306, -110 & 243, -110 & 126, -110 & 218, -110 & 110, -110 & 229, -285 \\
b$\delta$ ink & 172, 3 &  & 13, 38 &  &  & \textbf{2, 1} &  & -6, 49 &  &  &  &  &  & \textbf{13, 2} &  &  & 135, 176 \\
\hline si\rule[-0.5em]{0pt}{1.9em} & {\bfseries\itshape \char"00C6 \char"00C6\cgj\char"0307 \char"00C6\cgj\char"0331} & {\bfseries\itshape \char"00D0 \char"00D0\cgj\char"0307 \char"00D0\cgj\char"0331} & {\bfseries\itshape \char"00D8 \char"00D8\cgj\char"0307 \char"00D8\cgj\char"0331} & {\bfseries\itshape \char"00DE \char"00DE\cgj\char"0307 \char"00DE\cgj\char"0331} & {\bfseries\itshape \char"00DF \char"00DF\cgj\char"0307 \char"00DF\cgj\char"0331} & {\bfseries\itshape \char"00E6 \char"00E6\cgj\char"0307 \char"00E6\cgj\char"0331} & {\bfseries\itshape \char"00F0 \char"00F0\cgj\char"0307 \char"00F0\cgj\char"0331} & {\bfseries\itshape \char"00F8 \char"00F8\cgj\char"0307 \char"00F8\cgj\char"0331} & {\bfseries\itshape \char"00FE \char"00FE\cgj\char"0307 \char"00FE\cgj\char"0331} & {\bfseries\itshape \char"0110 \char"0110\cgj\char"0307 \char"0110\cgj\char"0331} & {\bfseries\itshape \char"0111 \char"0111\cgj\char"0307 \char"0111\cgj\char"0331} & {\bfseries\itshape \char"0126 \char"0126\cgj\char"0307 \char"0126\cgj\char"0331} & {\bfseries\itshape \char"0127 \char"0127\cgj\char"0307 \char"0127\cgj\char"0331} & {\bfseries\itshape \char"0131 \char"0131\cgj\char"0307 \char"0131\cgj\char"0331} & {\bfseries\itshape \char"0141 \char"0141\cgj\char"0307 \char"0141\cgj\char"0331} & {\bfseries\itshape \char"0142 \char"0142\cgj\char"0307 \char"0142\cgj\char"0331} & {\bfseries\itshape \char"014A \char"014A\cgj\char"0307 \char"014A\cgj\char"0331} \\
patch\rule[-0.5em]{0pt}{1.9em} & {\LibertinusSerifPatch\bfseries\itshape \char"00C6 \char"00C6\cgj\char"0307 \char"00C6\cgj\char"0331} & {\LibertinusSerifPatch\bfseries\itshape \char"00D0 \char"00D0\cgj\char"0307 \char"00D0\cgj\char"0331} & {\LibertinusSerifPatch\bfseries\itshape \char"00D8 \char"00D8\cgj\char"0307 \char"00D8\cgj\char"0331} & {\LibertinusSerifPatch\bfseries\itshape \char"00DE \char"00DE\cgj\char"0307 \char"00DE\cgj\char"0331} & {\LibertinusSerifPatch\bfseries\itshape \char"00DF \char"00DF\cgj\char"0307 \char"00DF\cgj\char"0331} & {\LibertinusSerifPatch\bfseries\itshape \char"00E6 \char"00E6\cgj\char"0307 \char"00E6\cgj\char"0331} & {\LibertinusSerifPatch\bfseries\itshape \char"00F0 \char"00F0\cgj\char"0307 \char"00F0\cgj\char"0331} & {\LibertinusSerifPatch\bfseries\itshape \char"00F8 \char"00F8\cgj\char"0307 \char"00F8\cgj\char"0331} & {\LibertinusSerifPatch\bfseries\itshape \char"00FE \char"00FE\cgj\char"0307 \char"00FE\cgj\char"0331} & {\LibertinusSerifPatch\bfseries\itshape \char"0110 \char"0110\cgj\char"0307 \char"0110\cgj\char"0331} & {\LibertinusSerifPatch\bfseries\itshape \char"0111 \char"0111\cgj\char"0307 \char"0111\cgj\char"0331} & {\LibertinusSerifPatch\bfseries\itshape \char"0126 \char"0126\cgj\char"0307 \char"0126\cgj\char"0331} & {\LibertinusSerifPatch\bfseries\itshape \char"0127 \char"0127\cgj\char"0307 \char"0127\cgj\char"0331} & {\LibertinusSerifPatch\bfseries\itshape \char"0131 \char"0131\cgj\char"0307 \char"0131\cgj\char"0331} & {\LibertinusSerifPatch\bfseries\itshape \char"0141 \char"0141\cgj\char"0307 \char"0141\cgj\char"0331} & {\LibertinusSerifPatch\bfseries\itshape \char"0142 \char"0142\cgj\char"0307 \char"0142\cgj\char"0331} & {\LibertinusSerifPatch\bfseries\itshape \char"014A \char"014A\cgj\char"0307 \char"014A\cgj\char"0331} \\
a & 500, 850 & 341, 850 & 351, 850 & 281, 850 &  & 350, 645 &  &  & 278, 646 &  &  &  &  & 132, 645 &  &  & 359, 804 \\
a? adv & 459, 850 & 367, 850 & 390, 850 & 303, 850 & 325, 890 & 339, 645 & 291, 909 & 235, 645 & 300, 890 & 367, 850 & 309, 890 & 396, 850 & 331, 890 & 144, 645 & 300, 850 & 190, 890 & 390, 850 \\
a$\delta$ adv & \textbf{41, 0} & -26, 0 & \textbf{-39, 0} & -22, 0 &  & 11, 0 &  &  & -22, -244 &  &  &  &  & \textbf{-12, 0} &  &  & -31, -46 \\
a? ink & 448, 850 & 347, 850 & 393, 850 & 285, 850 & 271, 890 & 350, 645 & 298, 909 & 236, 645 & 287, 890 & 333, 850 & 318, 890 & 415, 850 & 331, 890 & 157, 645 & 322, 850 & 207, 890 & 375, 850 \\
a$\delta$ ink & 52, 0 & \textbf{-6, 0} & -42, 0 & \textbf{-4, 0} &  & \textbf{0, 0} &  &  & \textbf{-9, -244} &  &  &  &  & -25, 0 &  &  & \textbf{-16, -46} \\
b & 549, -107 &  & 350, -110 &  &  & 327, -109 &  & 243, -110 &  &  &  &  &  & 139, -108 &  &  & 373, -109 \\
b? adv & 426, -110 & 334, -110 & 349, -148 & 270, -110 & 241, -319 & 347, -110 & 246, -110 & 237, -138 & 216, -319 & 334, -110 & 268, -110 & 363, -110 & 290, -110 & 153, -110 & 267, -110 & 148, -110 & 325, -268 \\
b$\delta$ adv & \textbf{123, 3} &  & \textbf{1, 38} &  &  & \textbf{-20, 1} &  & 6, 28 &  &  &  &  &  & \textbf{-14, 2} &  &  & \textbf{48, 159} \\
b? ink & 415, -110 & 314, -110 & 352, -148 & 252, -110 & 187, -319 & 359, -110 & 253, -110 & 239, -138 & 203, -319 & 300, -110 & 277, -110 & 382, -110 & 290, -110 & 166, -110 & 289, -110 & 166, -110 & 309, -268 \\
b$\delta$ ink & 134, 3 &  & -2, 38 &  &  & -32, 1 &  & \textbf{4, 28} &  &  &  &  &  & -27, 2 &  &  & 64, 159 \\

\end{tabular}
}
\end{table}
\clearpage


\begin{table}[p]
\captionsetup{justification=raggedright, singlelinecheck=false}
\caption{Lane comparison for Latin and IPA 2/8}
\label{table:lanes_latinIPA_2}
\setlength{\tabcolsep}{3.5pt}

{\footnotesize
\begin{tabular}{lrrrrrrrrrrrrrrrrr}
hex & 014B & 0152 & 0153 & 017F & 0181 & 0182 & 0184 & 0186 & 0187 & 0189 & 018A & 018B & 018E & 018F & 0190 & 0191 & 0192 \\
cat & Ll & Lu & Ll & Ll & Lu & Lu & Lu & Lu & Lu & Lu & Lu & Lu & Lu & Lu & Lu & Lu & Ll \\
\hline it\rule[-0.5em]{0pt}{1.9em} & {\itshape \char"014B \char"014B\cgj\char"0307 \char"014B\cgj\char"0331} & {\itshape \char"0152 \char"0152\cgj\char"0307 \char"0152\cgj\char"0331} & {\itshape \char"0153 \char"0153\cgj\char"0307 \char"0153\cgj\char"0331} & {\itshape \char"017F \char"017F\cgj\char"0307 \char"017F\cgj\char"0331} & {\itshape \char"0181 \char"0181\cgj\char"0307 \char"0181\cgj\char"0331} & {\itshape \char"0182 \char"0182\cgj\char"0307 \char"0182\cgj\char"0331} & {\itshape \char"0184 \char"0184\cgj\char"0307 \char"0184\cgj\char"0331} & {\itshape \char"0186 \char"0186\cgj\char"0307 \char"0186\cgj\char"0331} & {\itshape \char"0187 \char"0187\cgj\char"0307 \char"0187\cgj\char"0331} & {\itshape \char"0189 \char"0189\cgj\char"0307 \char"0189\cgj\char"0331} & {\itshape \char"018A \char"018A\cgj\char"0307 \char"018A\cgj\char"0331} & {\itshape \char"018B \char"018B\cgj\char"0307 \char"018B\cgj\char"0331} & {\itshape \char"018E \char"018E\cgj\char"0307 \char"018E\cgj\char"0331} & {\itshape \char"018F \char"018F\cgj\char"0307 \char"018F\cgj\char"0331} & {\itshape \char"0190 \char"0190\cgj\char"0307 \char"0190\cgj\char"0331} & {\itshape \char"0191 \char"0191\cgj\char"0307 \char"0191\cgj\char"0331} & {\itshape \char"0192 \char"0192\cgj\char"0307 \char"0192\cgj\char"0331} \\
patch\rule[-0.5em]{0pt}{1.9em} & {\LibertinusSerifPatch\itshape \char"014B \char"014B\cgj\char"0307 \char"014B\cgj\char"0331} & {\LibertinusSerifPatch\itshape \char"0152 \char"0152\cgj\char"0307 \char"0152\cgj\char"0331} & {\LibertinusSerifPatch\itshape \char"0153 \char"0153\cgj\char"0307 \char"0153\cgj\char"0331} & {\LibertinusSerifPatch\itshape \char"017F \char"017F\cgj\char"0307 \char"017F\cgj\char"0331} & {\LibertinusSerifPatch\itshape \char"0181 \char"0181\cgj\char"0307 \char"0181\cgj\char"0331} & {\LibertinusSerifPatch\itshape \char"0182 \char"0182\cgj\char"0307 \char"0182\cgj\char"0331} & {\LibertinusSerifPatch\itshape \char"0184 \char"0184\cgj\char"0307 \char"0184\cgj\char"0331} & {\LibertinusSerifPatch\itshape \char"0186 \char"0186\cgj\char"0307 \char"0186\cgj\char"0331} & {\LibertinusSerifPatch\itshape \char"0187 \char"0187\cgj\char"0307 \char"0187\cgj\char"0331} & {\LibertinusSerifPatch\itshape \char"0189 \char"0189\cgj\char"0307 \char"0189\cgj\char"0331} & {\LibertinusSerifPatch\itshape \char"018A \char"018A\cgj\char"0307 \char"018A\cgj\char"0331} & {\LibertinusSerifPatch\itshape \char"018B \char"018B\cgj\char"0307 \char"018B\cgj\char"0331} & {\LibertinusSerifPatch\itshape \char"018E \char"018E\cgj\char"0307 \char"018E\cgj\char"0331} & {\LibertinusSerifPatch\itshape \char"018F \char"018F\cgj\char"0307 \char"018F\cgj\char"0331} & {\LibertinusSerifPatch\itshape \char"0190 \char"0190\cgj\char"0307 \char"0190\cgj\char"0331} & {\LibertinusSerifPatch\itshape \char"0191 \char"0191\cgj\char"0307 \char"0191\cgj\char"0331} & {\LibertinusSerifPatch\itshape \char"0192 \char"0192\cgj\char"0307 \char"0192\cgj\char"0331} \\
a & 264, 646 & 575, 850 & 480, 645 &  & 327, 805 &  & 267, 805 &  &  &  & 306, 805 &  &  &  &  &  &  \\
a? adv & 382, 645 & 587, 850 & 490, 645 & 339, 890 & 488, 850 & 465, 850 & 455, 850 & 503, 850 & 503, 850 & 534, 850 & 553, 850 & 450, 850 & 465, 850 & 515, 850 & 439, 850 & 419, 850 & 414, 890 \\
a$\delta$ adv & -118, 1 & -12, 0 & -10, 0 &  & -161, -45 &  & -188, -45 &  &  &  & -247, -45 &  &  &  &  &  &  \\
a? ink & 368, 645 & 585, 850 & 489, 645 & 360, 890 & 466, 850 & 454, 850 & 445, 850 & 512, 850 & 560, 850 & 520, 850 & 527, 850 & 459, 850 & 478, 850 & 527, 850 & 431, 850 & 364, 850 & 402, 890 \\
a$\delta$ ink & \textbf{-104, 1} & \textbf{-10, 0} & \textbf{-9, 0} &  & \textbf{-139, -45} &  & \textbf{-178, -45} &  &  &  & \textbf{-221, -45} &  &  &  &  &  &  \\
b & 265, -111 &  &  &  & 328, -108 &  & 268, -108 &  &  &  & 369, -107 &  &  &  &  &  &  \\
b? adv & 177, -319 & 383, -110 & 329, -110 & 82, -319 & 284, -110 & 261, -110 & 251, -110 & 299, -110 & 299, -110 & 330, -110 & 349, -110 & 246, -110 & 261, -110 & 311, -110 & 235, -110 & 184, -257 & 170, -257 \\
b$\delta$ adv & \textbf{88, 208} &  &  &  & \textbf{44, 2} &  & \textbf{17, 2} &  &  &  & \textbf{20, 3} &  &  &  &  &  &  \\
b? ink & 163, -319 & 381, -110 & 329, -110 & 103, -319 & 262, -110 & 250, -110 & 241, -110 & 308, -110 & 356, -110 & 316, -110 & 323, -110 & 255, -110 & 274, -110 & 323, -110 & 227, -110 & 128, -257 & 158, -257 \\
b$\delta$ ink & 102, 208 &  &  &  & 66, 2 &  & 27, 2 &  &  &  & 46, 3 &  &  &  &  &  &  \\
\hline si\rule[-0.5em]{0pt}{1.9em} & {\bfseries\itshape \char"014B \char"014B\cgj\char"0307 \char"014B\cgj\char"0331} & {\bfseries\itshape \char"0152 \char"0152\cgj\char"0307 \char"0152\cgj\char"0331} & {\bfseries\itshape \char"0153 \char"0153\cgj\char"0307 \char"0153\cgj\char"0331} & {\bfseries\itshape \char"017F \char"017F\cgj\char"0307 \char"017F\cgj\char"0331} & {\bfseries\itshape \char"0181 \char"0181\cgj\char"0307 \char"0181\cgj\char"0331} & {\bfseries\itshape \char"0182 \char"0182\cgj\char"0307 \char"0182\cgj\char"0331} & {\bfseries\itshape \char"0184 \char"0184\cgj\char"0307 \char"0184\cgj\char"0331} & {\bfseries\itshape \char"0186 \char"0186\cgj\char"0307 \char"0186\cgj\char"0331} & {\bfseries\itshape \char"0187 \char"0187\cgj\char"0307 \char"0187\cgj\char"0331} & {\bfseries\itshape \char"0189 \char"0189\cgj\char"0307 \char"0189\cgj\char"0331} & {\bfseries\itshape \char"018A \char"018A\cgj\char"0307 \char"018A\cgj\char"0331} & {\bfseries\itshape \char"018B \char"018B\cgj\char"0307 \char"018B\cgj\char"0331} & {\bfseries\itshape \char"018E \char"018E\cgj\char"0307 \char"018E\cgj\char"0331} & {\bfseries\itshape \char"018F \char"018F\cgj\char"0307 \char"018F\cgj\char"0331} & {\bfseries\itshape \char"0190 \char"0190\cgj\char"0307 \char"0190\cgj\char"0331} & {\bfseries\itshape \char"0191 \char"0191\cgj\char"0307 \char"0191\cgj\char"0331} & {\bfseries\itshape \char"0192 \char"0192\cgj\char"0307 \char"0192\cgj\char"0331} \\
patch\rule[-0.5em]{0pt}{1.9em} & {\LibertinusSerifPatch\bfseries\itshape \char"014B \char"014B\cgj\char"0307 \char"014B\cgj\char"0331} & {\LibertinusSerifPatch\bfseries\itshape \char"0152 \char"0152\cgj\char"0307 \char"0152\cgj\char"0331} & {\LibertinusSerifPatch\bfseries\itshape \char"0153 \char"0153\cgj\char"0307 \char"0153\cgj\char"0331} & {\LibertinusSerifPatch\bfseries\itshape \char"017F \char"017F\cgj\char"0307 \char"017F\cgj\char"0331} & {\LibertinusSerifPatch\bfseries\itshape \char"0181 \char"0181\cgj\char"0307 \char"0181\cgj\char"0331} & {\LibertinusSerifPatch\bfseries\itshape \char"0182 \char"0182\cgj\char"0307 \char"0182\cgj\char"0331} & {\LibertinusSerifPatch\bfseries\itshape \char"0184 \char"0184\cgj\char"0307 \char"0184\cgj\char"0331} & {\LibertinusSerifPatch\bfseries\itshape \char"0186 \char"0186\cgj\char"0307 \char"0186\cgj\char"0331} & {\LibertinusSerifPatch\bfseries\itshape \char"0187 \char"0187\cgj\char"0307 \char"0187\cgj\char"0331} & {\LibertinusSerifPatch\bfseries\itshape \char"0189 \char"0189\cgj\char"0307 \char"0189\cgj\char"0331} & {\LibertinusSerifPatch\bfseries\itshape \char"018A \char"018A\cgj\char"0307 \char"018A\cgj\char"0331} & {\LibertinusSerifPatch\bfseries\itshape \char"018B \char"018B\cgj\char"0307 \char"018B\cgj\char"0331} & {\LibertinusSerifPatch\bfseries\itshape \char"018E \char"018E\cgj\char"0307 \char"018E\cgj\char"0331} & {\LibertinusSerifPatch\bfseries\itshape \char"018F \char"018F\cgj\char"0307 \char"018F\cgj\char"0331} & {\LibertinusSerifPatch\bfseries\itshape \char"0190 \char"0190\cgj\char"0307 \char"0190\cgj\char"0331} & {\LibertinusSerifPatch\bfseries\itshape \char"0191 \char"0191\cgj\char"0307 \char"0191\cgj\char"0331} & {\LibertinusSerifPatch\bfseries\itshape \char"0192 \char"0192\cgj\char"0307 \char"0192\cgj\char"0331} \\
a & 264, 646 & 480, 850 & 415, 645 &  & 327, 805 &  & 267, 805 &  &  &  & 306, 805 &  &  &  &  &  &  \\
a? adv & 253, 645 & 450, 850 & 384, 645 & 209, 903 & 359, 850 & 303, 850 & 306, 850 & 357, 850 & 359, 850 & 388, 850 & 424, 850 & 304, 850 & 320, 850 & 369, 850 & 293, 850 & 288, 850 & 269, 890 \\
a$\delta$ adv & \textbf{11, 1} & 30, 0 & 31, 0 &  & -32, -45 &  & -39, -45 &  &  &  & -118, -45 &  &  &  &  &  &  \\
a? ink & 236, 645 & 464, 850 & 392, 645 & 222, 903 & 329, 850 & 279, 850 & 289, 850 & 337, 850 & 398, 850 & 382, 850 & 393, 850 & 313, 850 & 333, 850 & 368, 850 & 287, 850 & 254, 850 & 259, 890 \\
a$\delta$ ink & 28, 1 & \textbf{16, 0} & \textbf{23, 0} &  & \textbf{-2, -45} &  & \textbf{-22, -45} &  &  &  & \textbf{-87, -45} &  &  &  &  &  &  \\
b & 265, -111 &  &  &  & 328, -108 &  & 268, -108 &  &  &  & 369, -107 &  &  &  &  &  &  \\
b? adv & 222, -301 & 417, -110 & 393, -110 & 121, -326 & 326, -110 & 270, -110 & 273, -110 & 324, -110 & 326, -110 & 355, -110 & 391, -110 & 271, -110 & 287, -110 & 336, -110 & 260, -110 & 220, -281 & 198, -255 \\
b$\delta$ adv & \textbf{43, 190} &  &  &  & \textbf{2, 2} &  & \textbf{-5, 2} &  &  &  & -22, 3 &  &  &  &  &  &  \\
b? ink & 206, -301 & 431, -110 & 401, -110 & 134, -326 & 296, -110 & 246, -110 & 256, -110 & 304, -110 & 365, -110 & 349, -110 & 360, -110 & 280, -110 & 300, -110 & 335, -110 & 254, -110 & 186, -281 & 188, -255 \\
b$\delta$ ink & 59, 190 &  &  &  & 32, 2 &  & 12, 2 &  &  &  & \textbf{9, 3} &  &  &  &  &  &  \\

\end{tabular}
}
\end{table}
\clearpage


\begin{table}[p]
\captionsetup{justification=raggedright, singlelinecheck=false}
\caption{Lane comparison for Latin and IPA 3/8}
\label{table:lanes_latinIPA_3}
\setlength{\tabcolsep}{3.5pt}

{\footnotesize
\begin{tabular}{lrrrrrrrrrrrrrrrrr}
hex & 0193 & 0194 & 0196 & 0197 & 0198 & 019C & 019D & 019F & 01A2 & 01A9 & 01B7 & 0237 & 024A & 0250 & 0251 & 0252 & 0253 \\
cat & Lu & Lu & Lu & Lu & Lu & Lu & Lu & Lu & Lu & Lu & Lu & Ll & Lu & Ll & Ll & Ll & Ll \\
\hline it\rule[-0.5em]{0pt}{1.9em} & {\itshape \char"0193 \char"0193\cgj\char"0307 \char"0193\cgj\char"0331} & {\itshape \char"0194 \char"0194\cgj\char"0307 \char"0194\cgj\char"0331} & {\itshape \char"0196 \char"0196\cgj\char"0307 \char"0196\cgj\char"0331} & {\itshape \char"0197 \char"0197\cgj\char"0307 \char"0197\cgj\char"0331} & {\itshape \char"0198 \char"0198\cgj\char"0307 \char"0198\cgj\char"0331} & {\itshape \char"019C \char"019C\cgj\char"0307 \char"019C\cgj\char"0331} & {\itshape \char"019D \char"019D\cgj\char"0307 \char"019D\cgj\char"0331} & {\itshape \char"019F \char"019F\cgj\char"0307 \char"019F\cgj\char"0331} & {\itshape \char"01A2 \char"01A2\cgj\char"0307 \char"01A2\cgj\char"0331} & {\itshape \char"01A9 \char"01A9\cgj\char"0307 \char"01A9\cgj\char"0331} & {\itshape \char"01B7 \char"01B7\cgj\char"0307 \char"01B7\cgj\char"0331} & {\itshape \char"0237 \char"0237\cgj\char"0307 \char"0237\cgj\char"0331} & {\itshape \char"024A \char"024A\cgj\char"0307 \char"024A\cgj\char"0331} & {\itshape \char"0250 \char"0250\cgj\char"0307 \char"0250\cgj\char"0331} & {\itshape \char"0251 \char"0251\cgj\char"0307 \char"0251\cgj\char"0331} & {\itshape \char"0252 \char"0252\cgj\char"0307 \char"0252\cgj\char"0331} & {\itshape \char"0253 \char"0253\cgj\char"0307 \char"0253\cgj\char"0331} \\
patch\rule[-0.5em]{0pt}{1.9em} & {\LibertinusSerifPatch\itshape \char"0193 \char"0193\cgj\char"0307 \char"0193\cgj\char"0331} & {\LibertinusSerifPatch\itshape \char"0194 \char"0194\cgj\char"0307 \char"0194\cgj\char"0331} & {\LibertinusSerifPatch\itshape \char"0196 \char"0196\cgj\char"0307 \char"0196\cgj\char"0331} & {\LibertinusSerifPatch\itshape \char"0197 \char"0197\cgj\char"0307 \char"0197\cgj\char"0331} & {\LibertinusSerifPatch\itshape \char"0198 \char"0198\cgj\char"0307 \char"0198\cgj\char"0331} & {\LibertinusSerifPatch\itshape \char"019C \char"019C\cgj\char"0307 \char"019C\cgj\char"0331} & {\LibertinusSerifPatch\itshape \char"019D \char"019D\cgj\char"0307 \char"019D\cgj\char"0331} & {\LibertinusSerifPatch\itshape \char"019F \char"019F\cgj\char"0307 \char"019F\cgj\char"0331} & {\LibertinusSerifPatch\itshape \char"01A2 \char"01A2\cgj\char"0307 \char"01A2\cgj\char"0331} & {\LibertinusSerifPatch\itshape \char"01A9 \char"01A9\cgj\char"0307 \char"01A9\cgj\char"0331} & {\LibertinusSerifPatch\itshape \char"01B7 \char"01B7\cgj\char"0307 \char"01B7\cgj\char"0331} & {\LibertinusSerifPatch\itshape \char"0237 \char"0237\cgj\char"0307 \char"0237\cgj\char"0331} & {\LibertinusSerifPatch\itshape \char"024A \char"024A\cgj\char"0307 \char"024A\cgj\char"0331} & {\LibertinusSerifPatch\itshape \char"0250 \char"0250\cgj\char"0307 \char"0250\cgj\char"0331} & {\LibertinusSerifPatch\itshape \char"0251 \char"0251\cgj\char"0307 \char"0251\cgj\char"0331} & {\LibertinusSerifPatch\itshape \char"0252 \char"0252\cgj\char"0307 \char"0252\cgj\char"0331} & {\LibertinusSerifPatch\itshape \char"0253 \char"0253\cgj\char"0307 \char"0253\cgj\char"0331} \\
a &  & 331, 805 &  &  &  & 431, 801 &  &  & 355, 804 &  &  & 140, 645 &  &  &  &  &  \\
a? adv & 541, 850 & 462, 850 & 322, 890 & 321, 850 & 499, 850 & 609, 850 & 528, 850 & 515, 850 & 622, 850 & 477, 850 & 485, 850 & 267, 645 &  & 365, 645 & 380, 645 & 391, 645 & 431, 896 \\
a$\delta$ adv &  & -131, -45 &  &  &  & -178, -49 &  &  & \textbf{-267, -46} &  &  & -127, 0 &  &  &  &  &  \\
a? ink & 537, 850 & 439, 850 & 327, 890 & 315, 850 & 494, 850 & 600, 850 & 496, 850 & 519, 850 & 632, 850 & 444, 850 & 485, 850 & 205, 645 &  & 361, 645 & 404, 645 & 379, 645 & 448, 896 \\
a$\delta$ ink &  & \textbf{-108, -45} &  &  &  & \textbf{-169, -49} &  &  & -277, -46 &  &  & \textbf{-65, 0} &  &  &  &  &  \\
b &  &  &  &  &  & 450, -108 &  &  & 350, -113 &  &  &  &  & 229, -114 & 212, -114 & 243, -114 & 212, -114 \\
b? adv & 337, -110 & 218, -297 & 110, -110 & 117, -110 & 295, -110 & 405, -110 & 290, -272 & 311, -110 & 393, -228 & 273, -110 & 241, -296 & 70, -281 &  & 205, -110 & 220, -110 & 231, -110 & 218, -110 \\
b$\delta$ adv &  &  &  &  &  & \textbf{45, 2} &  &  & \textbf{-43, 115} &  &  &  &  & \textbf{24, -4} & \textbf{-8, -4} & \textbf{12, -4} & \textbf{-6, -4} \\
b? ink & 333, -110 & 195, -297 & 115, -110 & 111, -110 & 290, -110 & 396, -110 & 257, -272 & 315, -110 & 403, -228 & 240, -110 & 241, -296 & 8, -281 &  & 201, -110 & 244, -110 & 219, -110 & 235, -110 \\
b$\delta$ ink &  &  &  &  &  & 54, 2 &  &  & -53, 115 &  &  &  &  & 28, -4 & -32, -4 & 24, -4 & -23, -4 \\
\hline si\rule[-0.5em]{0pt}{1.9em} & {\bfseries\itshape \char"0193 \char"0193\cgj\char"0307 \char"0193\cgj\char"0331} & {\bfseries\itshape \char"0194 \char"0194\cgj\char"0307 \char"0194\cgj\char"0331} & {\bfseries\itshape \char"0196 \char"0196\cgj\char"0307 \char"0196\cgj\char"0331} & {\bfseries\itshape \char"0197 \char"0197\cgj\char"0307 \char"0197\cgj\char"0331} & {\bfseries\itshape \char"0198 \char"0198\cgj\char"0307 \char"0198\cgj\char"0331} & {\bfseries\itshape \char"019C \char"019C\cgj\char"0307 \char"019C\cgj\char"0331} & {\bfseries\itshape \char"019D \char"019D\cgj\char"0307 \char"019D\cgj\char"0331} & {\bfseries\itshape \char"019F \char"019F\cgj\char"0307 \char"019F\cgj\char"0331} & {\bfseries\itshape \char"01A2 \char"01A2\cgj\char"0307 \char"01A2\cgj\char"0331} & {\bfseries\itshape \char"01A9 \char"01A9\cgj\char"0307 \char"01A9\cgj\char"0331} & {\bfseries\itshape \char"01B7 \char"01B7\cgj\char"0307 \char"01B7\cgj\char"0331} & {\bfseries\itshape \char"0237 \char"0237\cgj\char"0307 \char"0237\cgj\char"0331} & {\bfseries\itshape \char"024A \char"024A\cgj\char"0307 \char"024A\cgj\char"0331} & {\bfseries\itshape \char"0250 \char"0250\cgj\char"0307 \char"0250\cgj\char"0331} & {\bfseries\itshape \char"0251 \char"0251\cgj\char"0307 \char"0251\cgj\char"0331} & {\bfseries\itshape \char"0252 \char"0252\cgj\char"0307 \char"0252\cgj\char"0331} & {\bfseries\itshape \char"0253 \char"0253\cgj\char"0307 \char"0253\cgj\char"0331} \\
patch\rule[-0.5em]{0pt}{1.9em} & {\LibertinusSerifPatch\bfseries\itshape \char"0193 \char"0193\cgj\char"0307 \char"0193\cgj\char"0331} & {\LibertinusSerifPatch\bfseries\itshape \char"0194 \char"0194\cgj\char"0307 \char"0194\cgj\char"0331} & {\LibertinusSerifPatch\bfseries\itshape \char"0196 \char"0196\cgj\char"0307 \char"0196\cgj\char"0331} & {\LibertinusSerifPatch\bfseries\itshape \char"0197 \char"0197\cgj\char"0307 \char"0197\cgj\char"0331} & {\LibertinusSerifPatch\bfseries\itshape \char"0198 \char"0198\cgj\char"0307 \char"0198\cgj\char"0331} & {\LibertinusSerifPatch\bfseries\itshape \char"019C \char"019C\cgj\char"0307 \char"019C\cgj\char"0331} & {\LibertinusSerifPatch\bfseries\itshape \char"019D \char"019D\cgj\char"0307 \char"019D\cgj\char"0331} & {\LibertinusSerifPatch\bfseries\itshape \char"019F \char"019F\cgj\char"0307 \char"019F\cgj\char"0331} & {\LibertinusSerifPatch\bfseries\itshape \char"01A2 \char"01A2\cgj\char"0307 \char"01A2\cgj\char"0331} & {\LibertinusSerifPatch\bfseries\itshape \char"01A9 \char"01A9\cgj\char"0307 \char"01A9\cgj\char"0331} & {\LibertinusSerifPatch\bfseries\itshape \char"01B7 \char"01B7\cgj\char"0307 \char"01B7\cgj\char"0331} & {\LibertinusSerifPatch\bfseries\itshape \char"0237 \char"0237\cgj\char"0307 \char"0237\cgj\char"0331} & {\LibertinusSerifPatch\bfseries\itshape \char"024A \char"024A\cgj\char"0307 \char"024A\cgj\char"0331} & {\LibertinusSerifPatch\bfseries\itshape \char"0250 \char"0250\cgj\char"0307 \char"0250\cgj\char"0331} & {\LibertinusSerifPatch\bfseries\itshape \char"0251 \char"0251\cgj\char"0307 \char"0251\cgj\char"0331} & {\LibertinusSerifPatch\bfseries\itshape \char"0252 \char"0252\cgj\char"0307 \char"0252\cgj\char"0331} & {\LibertinusSerifPatch\bfseries\itshape \char"0253 \char"0253\cgj\char"0307 \char"0253\cgj\char"0331} \\
a &  & 331, 805 &  &  &  & 431, 801 &  &  & 359, 804 &  &  & 140, 645 &  &  & 228, 645 & 264, 645 &  \\
a? adv & 395, 850 & 316, 850 & 194, 890 & 181, 850 & 353, 850 & 463, 850 & 383, 850 & 387, 850 & 464, 850 & 319, 850 & 329, 850 & 121, 645 &  & 221, 645 & 247, 645 & 247, 645 & 284, 890 \\
a$\delta$ adv &  & \textbf{15, -45} &  &  &  & -32, -49 &  &  & \textbf{-105, -46} &  &  & \textbf{19, 0} &  &  & \textbf{-19, 0} & \textbf{17, 0} &  \\
a? ink & 396, 850 & 299, 850 & 193, 890 & 185, 850 & 352, 850 & 458, 850 & 353, 850 & 394, 850 & 474, 850 & 290, 850 & 326, 850 & 89, 645 &  & 203, 645 & 262, 645 & 239, 645 & 302, 890 \\
a$\delta$ ink &  & 32, -45 &  &  &  & \textbf{-27, -49} &  &  & -115, -46 &  &  & 51, 0 &  &  & -34, 0 & 25, 0 &  \\
b &  &  &  &  &  & 450, -108 &  &  & 354, -113 &  & 236, -319 &  &  & 229, -114 & 212, -114 & 259, -114 & 212, -114 \\
b? adv & 362, -110 & 245, -295 & 153, -110 & 148, -110 & 320, -110 & 430, -110 & 317, -270 & 354, -110 & 407, -227 & 286, -110 & 258, -293 & 93, -292 &  & 230, -110 & 256, -110 & 256, -110 & 243, -110 \\
b$\delta$ adv &  &  &  &  &  & \textbf{20, 2} &  &  & \textbf{-53, 114} &  & -22, -26 &  &  & \textbf{-1, -4} & \textbf{-44, -4} & \textbf{3, -4} & \textbf{-31, -4} \\
b? ink & 363, -110 & 228, -295 & 152, -110 & 152, -110 & 319, -110 & 425, -110 & 287, -270 & 361, -110 & 417, -227 & 257, -110 & 255, -293 & 61, -292 &  & 212, -110 & 271, -110 & 248, -110 & 261, -110 \\
b$\delta$ ink &  &  &  &  &  & 25, 2 &  &  & -63, 114 &  & \textbf{-19, -26} &  &  & 17, -4 & -59, -4 & 11, -4 & -49, -4 \\

\end{tabular}
}
\end{table}
\clearpage


\begin{table}[p]
\captionsetup{justification=raggedright, singlelinecheck=false}
\caption{Lane comparison for Latin and IPA 4/8}
\label{table:lanes_latinIPA_4}
\setlength{\tabcolsep}{3.5pt}

{\footnotesize
\begin{tabular}{lrrrrrrrrrrrrrrrrr}
hex & 0254 & 0255 & 0256 & 0257 & 0258 & 0259 & 025A & 025B & 025C & 025D & 025E & 025F & 0260 & 0261 & 0262 & 0263 & 0264 \\
cat & Ll & Ll & Ll & Ll & Ll & Ll & Ll & Ll & Ll & Ll & Ll & Ll & Ll & Ll & Ll & Ll & Ll \\
\hline it\rule[-0.5em]{0pt}{1.9em} & {\itshape \char"0254 \char"0254\cgj\char"0307 \char"0254\cgj\char"0331} & {\itshape \char"0255 \char"0255\cgj\char"0307 \char"0255\cgj\char"0331} & {\itshape \char"0256 \char"0256\cgj\char"0307 \char"0256\cgj\char"0331} & {\itshape \char"0257 \char"0257\cgj\char"0307 \char"0257\cgj\char"0331} & {\itshape \char"0258 \char"0258\cgj\char"0307 \char"0258\cgj\char"0331} & {\itshape \char"0259 \char"0259\cgj\char"0307 \char"0259\cgj\char"0331} & {\itshape \char"025A \char"025A\cgj\char"0307 \char"025A\cgj\char"0331} & {\itshape \char"025B \char"025B\cgj\char"0307 \char"025B\cgj\char"0331} & {\itshape \char"025C \char"025C\cgj\char"0307 \char"025C\cgj\char"0331} & {\itshape \char"025D \char"025D\cgj\char"0307 \char"025D\cgj\char"0331} & {\itshape \char"025E \char"025E\cgj\char"0307 \char"025E\cgj\char"0331} & {\itshape \char"025F \char"025F\cgj\char"0307 \char"025F\cgj\char"0331} & {\itshape \char"0260 \char"0260\cgj\char"0307 \char"0260\cgj\char"0331} & {\itshape \char"0261 \char"0261\cgj\char"0307 \char"0261\cgj\char"0331} & {\itshape \char"0262 \char"0262\cgj\char"0307 \char"0262\cgj\char"0331} & {\itshape \char"0263 \char"0263\cgj\char"0307 \char"0263\cgj\char"0331} & {\itshape \char"0264 \char"0264\cgj\char"0307 \char"0264\cgj\char"0331} \\
patch\rule[-0.5em]{0pt}{1.9em} & {\LibertinusSerifPatch\itshape \char"0254 \char"0254\cgj\char"0307 \char"0254\cgj\char"0331} & {\LibertinusSerifPatch\itshape \char"0255 \char"0255\cgj\char"0307 \char"0255\cgj\char"0331} & {\LibertinusSerifPatch\itshape \char"0256 \char"0256\cgj\char"0307 \char"0256\cgj\char"0331} & {\LibertinusSerifPatch\itshape \char"0257 \char"0257\cgj\char"0307 \char"0257\cgj\char"0331} & {\LibertinusSerifPatch\itshape \char"0258 \char"0258\cgj\char"0307 \char"0258\cgj\char"0331} & {\LibertinusSerifPatch\itshape \char"0259 \char"0259\cgj\char"0307 \char"0259\cgj\char"0331} & {\LibertinusSerifPatch\itshape \char"025A \char"025A\cgj\char"0307 \char"025A\cgj\char"0331} & {\LibertinusSerifPatch\itshape \char"025B \char"025B\cgj\char"0307 \char"025B\cgj\char"0331} & {\LibertinusSerifPatch\itshape \char"025C \char"025C\cgj\char"0307 \char"025C\cgj\char"0331} & {\LibertinusSerifPatch\itshape \char"025D \char"025D\cgj\char"0307 \char"025D\cgj\char"0331} & {\LibertinusSerifPatch\itshape \char"025E \char"025E\cgj\char"0307 \char"025E\cgj\char"0331} & {\LibertinusSerifPatch\itshape \char"025F \char"025F\cgj\char"0307 \char"025F\cgj\char"0331} & {\LibertinusSerifPatch\itshape \char"0260 \char"0260\cgj\char"0307 \char"0260\cgj\char"0331} & {\LibertinusSerifPatch\itshape \char"0261 \char"0261\cgj\char"0307 \char"0261\cgj\char"0331} & {\LibertinusSerifPatch\itshape \char"0262 \char"0262\cgj\char"0307 \char"0262\cgj\char"0331} & {\LibertinusSerifPatch\itshape \char"0263 \char"0263\cgj\char"0307 \char"0263\cgj\char"0331} & {\LibertinusSerifPatch\itshape \char"0264 \char"0264\cgj\char"0307 \char"0264\cgj\char"0331} \\
a & 322, 654 &  & 218, 645 & 218, 645 & 204, 645 &  &  & 322, 654 &  &  &  &  &  &  & 265, 645 &  &  \\
a? adv & 332, 645 & 352, 645 & 465, 899 & 464, 890 & 356, 645 & 356, 645 & 418, 687 & 338, 645 & 331, 645 & 413, 775 & 375, 645 & 305, 645 & 431, 896 & 378, 645 & 388, 645 & 387, 645 & 338, 645 \\
a$\delta$ adv & -10, 9 &  & \textbf{-247, -254} & \textbf{-246, -245} & \textbf{-152, 0} &  &  & \textbf{-16, 9} &  &  &  &  &  &  & -123, 0 &  &  \\
a? ink & 327, 645 & 355, 645 & 500, 899 & 502, 890 & 359, 645 & 358, 645 & 459, 687 & 340, 645 & 331, 645 & 425, 775 & 376, 645 & 273, 645 & 521, 896 & 366, 645 & 386, 645 & 393, 645 & 345, 645 \\
a$\delta$ ink & \textbf{-5, 9} &  & -282, -254 & -284, -245 & -155, 0 &  &  & -18, 9 &  &  &  &  &  &  & \textbf{-121, 0} &  &  \\
b & 172, -114 & 212, -114 & 212, -114 & 212, -114 & 212, -114 & 212, -114 & 212, -114 & 212, -114 & 184, -114 & 212, -114 & 256, -114 &  &  &  & 248, -114 &  & 177, -114 \\
b? adv & 171, -110 & 171, -208 & 215, -273 & 251, -110 & 195, -110 & 195, -110 & 248, -110 & 178, -110 & 170, -110 & 225, -110 & 209, -133 & 100, -319 & 180, -286 & 180, -286 & 227, -110 & 192, -272 & 178, -106 \\
b$\delta$ adv & \textbf{1, -4} & 41, 94 & \textbf{-3, 159} & \textbf{-39, -4} & 17, -4 & 17, -4 & \textbf{-36, -4} & 34, -4 & 14, -4 & \textbf{-13, -4} & 47, 19 &  &  &  & \textbf{21, -4} &  & \textbf{-1, -8} \\
b? ink & 167, -110 & 174, -208 & 251, -273 & 290, -110 & 199, -110 & 198, -110 & 290, -110 & 180, -110 & 171, -110 & 237, -110 & 211, -133 & 68, -319 & 270, -286 & 168, -286 & 226, -110 & 198, -272 & 185, -106 \\
b$\delta$ ink & 5, -4 & \textbf{38, 94} & -39, 159 & -78, -4 & \textbf{13, -4} & \textbf{14, -4} & -78, -4 & \textbf{32, -4} & \textbf{13, -4} & -25, -4 & \textbf{45, 19} &  &  &  & 22, -4 &  & -8, -8 \\
\hline si\rule[-0.5em]{0pt}{1.9em} & {\bfseries\itshape \char"0254 \char"0254\cgj\char"0307 \char"0254\cgj\char"0331} & {\bfseries\itshape \char"0255 \char"0255\cgj\char"0307 \char"0255\cgj\char"0331} & {\bfseries\itshape \char"0256 \char"0256\cgj\char"0307 \char"0256\cgj\char"0331} & {\bfseries\itshape \char"0257 \char"0257\cgj\char"0307 \char"0257\cgj\char"0331} & {\bfseries\itshape \char"0258 \char"0258\cgj\char"0307 \char"0258\cgj\char"0331} & {\bfseries\itshape \char"0259 \char"0259\cgj\char"0307 \char"0259\cgj\char"0331} & {\bfseries\itshape \char"025A \char"025A\cgj\char"0307 \char"025A\cgj\char"0331} & {\bfseries\itshape \char"025B \char"025B\cgj\char"0307 \char"025B\cgj\char"0331} & {\bfseries\itshape \char"025C \char"025C\cgj\char"0307 \char"025C\cgj\char"0331} & {\bfseries\itshape \char"025D \char"025D\cgj\char"0307 \char"025D\cgj\char"0331} & {\bfseries\itshape \char"025E \char"025E\cgj\char"0307 \char"025E\cgj\char"0331} & {\bfseries\itshape \char"025F \char"025F\cgj\char"0307 \char"025F\cgj\char"0331} & {\bfseries\itshape \char"0260 \char"0260\cgj\char"0307 \char"0260\cgj\char"0331} & {\bfseries\itshape \char"0261 \char"0261\cgj\char"0307 \char"0261\cgj\char"0331} & {\bfseries\itshape \char"0262 \char"0262\cgj\char"0307 \char"0262\cgj\char"0331} & {\bfseries\itshape \char"0263 \char"0263\cgj\char"0307 \char"0263\cgj\char"0331} & {\bfseries\itshape \char"0264 \char"0264\cgj\char"0307 \char"0264\cgj\char"0331} \\
patch\rule[-0.5em]{0pt}{1.9em} & {\LibertinusSerifPatch\bfseries\itshape \char"0254 \char"0254\cgj\char"0307 \char"0254\cgj\char"0331} & {\LibertinusSerifPatch\bfseries\itshape \char"0255 \char"0255\cgj\char"0307 \char"0255\cgj\char"0331} & {\LibertinusSerifPatch\bfseries\itshape \char"0256 \char"0256\cgj\char"0307 \char"0256\cgj\char"0331} & {\LibertinusSerifPatch\bfseries\itshape \char"0257 \char"0257\cgj\char"0307 \char"0257\cgj\char"0331} & {\LibertinusSerifPatch\bfseries\itshape \char"0258 \char"0258\cgj\char"0307 \char"0258\cgj\char"0331} & {\LibertinusSerifPatch\bfseries\itshape \char"0259 \char"0259\cgj\char"0307 \char"0259\cgj\char"0331} & {\LibertinusSerifPatch\bfseries\itshape \char"025A \char"025A\cgj\char"0307 \char"025A\cgj\char"0331} & {\LibertinusSerifPatch\bfseries\itshape \char"025B \char"025B\cgj\char"0307 \char"025B\cgj\char"0331} & {\LibertinusSerifPatch\bfseries\itshape \char"025C \char"025C\cgj\char"0307 \char"025C\cgj\char"0331} & {\LibertinusSerifPatch\bfseries\itshape \char"025D \char"025D\cgj\char"0307 \char"025D\cgj\char"0331} & {\LibertinusSerifPatch\bfseries\itshape \char"025E \char"025E\cgj\char"0307 \char"025E\cgj\char"0331} & {\LibertinusSerifPatch\bfseries\itshape \char"025F \char"025F\cgj\char"0307 \char"025F\cgj\char"0331} & {\LibertinusSerifPatch\bfseries\itshape \char"0260 \char"0260\cgj\char"0307 \char"0260\cgj\char"0331} & {\LibertinusSerifPatch\bfseries\itshape \char"0261 \char"0261\cgj\char"0307 \char"0261\cgj\char"0331} & {\LibertinusSerifPatch\bfseries\itshape \char"0262 \char"0262\cgj\char"0307 \char"0262\cgj\char"0331} & {\LibertinusSerifPatch\bfseries\itshape \char"0263 \char"0263\cgj\char"0307 \char"0263\cgj\char"0331} & {\LibertinusSerifPatch\bfseries\itshape \char"0264 \char"0264\cgj\char"0307 \char"0264\cgj\char"0331} \\
a & 160, 645 &  & 218, 645 & 218, 645 & 204, 645 &  &  & 209, 645 & 194, 645 &  &  &  &  &  & 265, 645 &  &  \\
a? adv & 199, 645 & 204, 645 & 302, 890 & 318, 890 & 207, 645 & 218, 645 & 299, 693 & 194, 645 & 194, 645 & 261, 762 & 232, 645 & 141, 645 & 309, 890 & 238, 645 & 244, 645 & 255, 645 & 185, 645 \\
a$\delta$ adv & -39, 0 &  & \textbf{-84, -245} & \textbf{-100, -245} & \textbf{-3, 0} &  &  & 15, 0 & \textbf{0, 0} &  &  &  &  &  & 21, 0 &  &  \\
a? ink & 184, 645 & 220, 645 & 355, 890 & 366, 890 & 215, 645 & 212, 645 & 319, 693 & 197, 645 & 199, 645 & 288, 762 & 237, 645 & 126, 645 & 378, 890 & 225, 645 & 252, 645 & 264, 645 & 195, 645 \\
a$\delta$ ink & \textbf{-24, 0} &  & -137, -245 & -148, -245 & -11, 0 &  &  & \textbf{12, 0} & -5, 0 &  &  &  &  &  & \textbf{13, 0} &  &  \\
b & 177, -114 & 212, -114 & 212, -114 & 212, -114 & 212, -114 &  & 212, -114 & 212, -114 & 184, -114 & 212, -114 & 256, -114 &  &  &  & 248, -110 &  & 177, -114 \\
b? adv & 207, -110 & 192, -209 & 227, -271 & 276, -110 & 216, -110 & 227, -110 & 298, -110 & 203, -110 & 203, -110 & 246, -110 & 241, -110 & 107, -319 & 232, -284 & 204, -319 & 252, -110 & 228, -285 & 194, -110 \\
b$\delta$ adv & -30, -4 & 20, 95 & \textbf{-15, 157} & \textbf{-64, -4} & \textbf{-4, -4} &  & \textbf{-86, -4} & 9, -4 & \textbf{-19, -4} & \textbf{-34, -4} & 15, -4 &  &  &  & \textbf{-4, 0} &  & \textbf{-17, -4} \\
b? ink & 193, -110 & 208, -209 & 281, -271 & 325, -110 & 224, -110 & 221, -110 & 318, -110 & 206, -110 & 208, -110 & 273, -110 & 246, -110 & 92, -319 & 301, -284 & 191, -319 & 261, -110 & 237, -285 & 204, -110 \\
b$\delta$ ink & \textbf{-16, -4} & \textbf{4, 95} & -69, 157 & -113, -4 & -12, -4 &  & -106, -4 & \textbf{6, -4} & -24, -4 & -61, -4 & \textbf{10, -4} &  &  &  & -13, 0 &  & -27, -4 \\

\end{tabular}
}
\end{table}
\clearpage


\begin{table}[p]
\captionsetup{justification=raggedright, singlelinecheck=false}
\caption{Lane comparison for Latin and IPA 5/8}
\label{table:lanes_latinIPA_5}
\setlength{\tabcolsep}{3.5pt}

{\footnotesize
\begin{tabular}{lrrrrrrrrrrrrrrrrr}
hex & 0265 & 0266 & 0267 & 0268 & 0269 & 026A & 026B & 026C & 026D & 026E & 026F & 0270 & 0271 & 0272 & 0273 & 0274 & 0275 \\
cat & Ll & Ll & Ll & Ll & Ll & Ll & Ll & Ll & Ll & Ll & Ll & Ll & Ll & Ll & Ll & Ll & Ll \\
\hline it\rule[-0.5em]{0pt}{1.9em} & {\itshape \char"0265 \char"0265\cgj\char"0307 \char"0265\cgj\char"0331} & {\itshape \char"0266 \char"0266\cgj\char"0307 \char"0266\cgj\char"0331} & {\itshape \char"0267 \char"0267\cgj\char"0307 \char"0267\cgj\char"0331} & {\itshape \char"0268 \char"0268\cgj\char"0307 \char"0268\cgj\char"0331} & {\itshape \char"0269 \char"0269\cgj\char"0307 \char"0269\cgj\char"0331} & {\itshape \char"026A \char"026A\cgj\char"0307 \char"026A\cgj\char"0331} & {\itshape \char"026B \char"026B\cgj\char"0307 \char"026B\cgj\char"0331} & {\itshape \char"026C \char"026C\cgj\char"0307 \char"026C\cgj\char"0331} & {\itshape \char"026D \char"026D\cgj\char"0307 \char"026D\cgj\char"0331} & {\itshape \char"026E \char"026E\cgj\char"0307 \char"026E\cgj\char"0331} & {\itshape \char"026F \char"026F\cgj\char"0307 \char"026F\cgj\char"0331} & {\itshape \char"0270 \char"0270\cgj\char"0307 \char"0270\cgj\char"0331} & {\itshape \char"0271 \char"0271\cgj\char"0307 \char"0271\cgj\char"0331} & {\itshape \char"0272 \char"0272\cgj\char"0307 \char"0272\cgj\char"0331} & {\itshape \char"0273 \char"0273\cgj\char"0307 \char"0273\cgj\char"0331} & {\itshape \char"0274 \char"0274\cgj\char"0307 \char"0274\cgj\char"0331} & {\itshape \char"0275 \char"0275\cgj\char"0307 \char"0275\cgj\char"0331} \\
patch\rule[-0.5em]{0pt}{1.9em} & {\LibertinusSerifPatch\itshape \char"0265 \char"0265\cgj\char"0307 \char"0265\cgj\char"0331} & {\LibertinusSerifPatch\itshape \char"0266 \char"0266\cgj\char"0307 \char"0266\cgj\char"0331} & {\LibertinusSerifPatch\itshape \char"0267 \char"0267\cgj\char"0307 \char"0267\cgj\char"0331} & {\LibertinusSerifPatch\itshape \char"0268 \char"0268\cgj\char"0307 \char"0268\cgj\char"0331} & {\LibertinusSerifPatch\itshape \char"0269 \char"0269\cgj\char"0307 \char"0269\cgj\char"0331} & {\LibertinusSerifPatch\itshape \char"026A \char"026A\cgj\char"0307 \char"026A\cgj\char"0331} & {\LibertinusSerifPatch\itshape \char"026B \char"026B\cgj\char"0307 \char"026B\cgj\char"0331} & {\LibertinusSerifPatch\itshape \char"026C \char"026C\cgj\char"0307 \char"026C\cgj\char"0331} & {\LibertinusSerifPatch\itshape \char"026D \char"026D\cgj\char"0307 \char"026D\cgj\char"0331} & {\LibertinusSerifPatch\itshape \char"026E \char"026E\cgj\char"0307 \char"026E\cgj\char"0331} & {\LibertinusSerifPatch\itshape \char"026F \char"026F\cgj\char"0307 \char"026F\cgj\char"0331} & {\LibertinusSerifPatch\itshape \char"0270 \char"0270\cgj\char"0307 \char"0270\cgj\char"0331} & {\LibertinusSerifPatch\itshape \char"0271 \char"0271\cgj\char"0307 \char"0271\cgj\char"0331} & {\LibertinusSerifPatch\itshape \char"0272 \char"0272\cgj\char"0307 \char"0272\cgj\char"0331} & {\LibertinusSerifPatch\itshape \char"0273 \char"0273\cgj\char"0307 \char"0273\cgj\char"0331} & {\LibertinusSerifPatch\itshape \char"0274 \char"0274\cgj\char"0307 \char"0274\cgj\char"0331} & {\LibertinusSerifPatch\itshape \char"0275 \char"0275\cgj\char"0307 \char"0275\cgj\char"0331} \\
a & 258, 645 & 299, 565 &  &  & 123, 645 & 128, 645 &  &  &  &  & 407, 645 &  & 383, 645 & 264, 646 & 263, 645 &  &  \\
a? adv & 396, 645 & 449, 890 & 449, 890 & 316, 810 & 266, 645 & 281, 645 & 333, 890 & 368, 890 & 341, 890 & 447, 890 & 517, 645 & 515, 645 & 515, 645 & 396, 645 & 367, 645 & 418, 645 & 379, 645 \\
a$\delta$ adv & \textbf{-138, 0} & \textbf{-150, -325} &  &  & \textbf{-143, 0} & -153, 0 &  &  &  &  & \textbf{-110, 0} &  & -132, 0 & -132, 1 & \textbf{-104, 0} &  &  \\
a? ink & 402, 645 & 487, 890 & 460, 890 & 324, 810 & 297, 645 & 277, 645 & 331, 890 & 367, 890 & 354, 890 & 445, 890 & 521, 645 & 526, 645 & 489, 645 & 362, 645 & 416, 645 & 420, 645 & 381, 645 \\
a$\delta$ ink & -144, 0 & -188, -325 &  &  & -174, 0 & \textbf{-149, 0} &  &  &  &  & -114, 0 &  & \textbf{-106, 0} & \textbf{-98, 1} & -153, 0 &  &  \\
b & 219, -129 & 240, -114 &  & 134, -114 & 153, -114 & 139, -114 & 146, -114 & 218, -120 &  &  & 410, -114 & 294, -114 & 383, -105 & 265, -111 & 262, -112 & 247, -114 &  \\
b? adv & 191, -319 & 236, -110 & 192, -319 & 121, -110 & 106, -110 & 120, -110 & 120, -110 & 156, -110 & 91, -287 & 192, -309 & 354, -123 & 310, -319 & 310, -319 & 191, -319 & 162, -319 & 258, -110 & 218, -110 \\
b$\delta$ adv & 28, 190 & \textbf{4, -4} &  & 13, -4 & 47, -4 & \textbf{19, -4} & \textbf{26, -4} & \textbf{62, -10} &  &  & 56, 9 & \textbf{-16, 205} & \textbf{73, 214} & \textbf{74, 208} & 100, 207 & \textbf{-11, -4} &  \\
b? ink & 197, -319 & 275, -110 & 203, -319 & 129, -110 & 137, -110 & 117, -110 & 119, -110 & 155, -110 & 104, -287 & 190, -309 & 358, -123 & 321, -319 & 284, -319 & 157, -319 & 211, -319 & 260, -110 & 221, -110 \\
b$\delta$ ink & \textbf{22, 190} & -35, -4 &  & \textbf{5, -4} & \textbf{16, -4} & 22, -4 & 27, -4 & 63, -10 &  &  & \textbf{52, 9} & -27, 205 & 99, 214 & 108, 208 & \textbf{51, 207} & -13, -4 &  \\
\hline si\rule[-0.5em]{0pt}{1.9em} & {\bfseries\itshape \char"0265 \char"0265\cgj\char"0307 \char"0265\cgj\char"0331} & {\bfseries\itshape \char"0266 \char"0266\cgj\char"0307 \char"0266\cgj\char"0331} & {\bfseries\itshape \char"0267 \char"0267\cgj\char"0307 \char"0267\cgj\char"0331} & {\bfseries\itshape \char"0268 \char"0268\cgj\char"0307 \char"0268\cgj\char"0331} & {\bfseries\itshape \char"0269 \char"0269\cgj\char"0307 \char"0269\cgj\char"0331} & {\bfseries\itshape \char"026A \char"026A\cgj\char"0307 \char"026A\cgj\char"0331} & {\bfseries\itshape \char"026B \char"026B\cgj\char"0307 \char"026B\cgj\char"0331} & {\bfseries\itshape \char"026C \char"026C\cgj\char"0307 \char"026C\cgj\char"0331} & {\bfseries\itshape \char"026D \char"026D\cgj\char"0307 \char"026D\cgj\char"0331} & {\bfseries\itshape \char"026E \char"026E\cgj\char"0307 \char"026E\cgj\char"0331} & {\bfseries\itshape \char"026F \char"026F\cgj\char"0307 \char"026F\cgj\char"0331} & {\bfseries\itshape \char"0270 \char"0270\cgj\char"0307 \char"0270\cgj\char"0331} & {\bfseries\itshape \char"0271 \char"0271\cgj\char"0307 \char"0271\cgj\char"0331} & {\bfseries\itshape \char"0272 \char"0272\cgj\char"0307 \char"0272\cgj\char"0331} & {\bfseries\itshape \char"0273 \char"0273\cgj\char"0307 \char"0273\cgj\char"0331} & {\bfseries\itshape \char"0274 \char"0274\cgj\char"0307 \char"0274\cgj\char"0331} & {\bfseries\itshape \char"0275 \char"0275\cgj\char"0307 \char"0275\cgj\char"0331} \\
patch\rule[-0.5em]{0pt}{1.9em} & {\LibertinusSerifPatch\bfseries\itshape \char"0265 \char"0265\cgj\char"0307 \char"0265\cgj\char"0331} & {\LibertinusSerifPatch\bfseries\itshape \char"0266 \char"0266\cgj\char"0307 \char"0266\cgj\char"0331} & {\LibertinusSerifPatch\bfseries\itshape \char"0267 \char"0267\cgj\char"0307 \char"0267\cgj\char"0331} & {\LibertinusSerifPatch\bfseries\itshape \char"0268 \char"0268\cgj\char"0307 \char"0268\cgj\char"0331} & {\LibertinusSerifPatch\bfseries\itshape \char"0269 \char"0269\cgj\char"0307 \char"0269\cgj\char"0331} & {\LibertinusSerifPatch\bfseries\itshape \char"026A \char"026A\cgj\char"0307 \char"026A\cgj\char"0331} & {\LibertinusSerifPatch\bfseries\itshape \char"026B \char"026B\cgj\char"0307 \char"026B\cgj\char"0331} & {\LibertinusSerifPatch\bfseries\itshape \char"026C \char"026C\cgj\char"0307 \char"026C\cgj\char"0331} & {\LibertinusSerifPatch\bfseries\itshape \char"026D \char"026D\cgj\char"0307 \char"026D\cgj\char"0331} & {\LibertinusSerifPatch\bfseries\itshape \char"026E \char"026E\cgj\char"0307 \char"026E\cgj\char"0331} & {\LibertinusSerifPatch\bfseries\itshape \char"026F \char"026F\cgj\char"0307 \char"026F\cgj\char"0331} & {\LibertinusSerifPatch\bfseries\itshape \char"0270 \char"0270\cgj\char"0307 \char"0270\cgj\char"0331} & {\LibertinusSerifPatch\bfseries\itshape \char"0271 \char"0271\cgj\char"0307 \char"0271\cgj\char"0331} & {\LibertinusSerifPatch\bfseries\itshape \char"0272 \char"0272\cgj\char"0307 \char"0272\cgj\char"0331} & {\LibertinusSerifPatch\bfseries\itshape \char"0273 \char"0273\cgj\char"0307 \char"0273\cgj\char"0331} & {\LibertinusSerifPatch\bfseries\itshape \char"0274 \char"0274\cgj\char"0307 \char"0274\cgj\char"0331} & {\LibertinusSerifPatch\bfseries\itshape \char"0275 \char"0275\cgj\char"0307 \char"0275\cgj\char"0331} \\
a & 258, 645 & 299, 565 & 299, 565 & 137, 795 & 123, 645 & 128, 645 &  &  &  &  & 407, 645 &  & 383, 645 & 264, 646 & 263, 645 &  &  \\
a? adv & 258, 645 & 302, 890 & 301, 890 & 180, 823 & 138, 645 & 137, 645 & 187, 890 & 223, 890 & 187, 890 & 301, 890 & 372, 640 & 373, 645 & 373, 645 & 254, 645 & 254, 645 & 275, 645 & 229, 645 \\
a$\delta$ adv & \textbf{0, 0} & \textbf{-3, -325} & \textbf{-2, -325} & \textbf{-43, -28} & \textbf{-15, 0} & -9, 0 &  &  &  &  & \textbf{35, 5} &  & \textbf{10, 0} & \textbf{10, 1} & \textbf{9, 0} &  &  \\
a? ink & 267, 645 & 305, 890 & 278, 890 & 191, 823 & 152, 645 & 136, 645 & 220, 890 & 234, 890 & 243, 890 & 299, 890 & 368, 640 & 381, 645 & 348, 645 & 224, 645 & 285, 645 & 277, 645 & 258, 645 \\
a$\delta$ ink & -9, 0 & -6, -325 & 21, -325 & -54, -28 & -29, 0 & \textbf{-8, 0} &  &  &  &  & 39, 5 &  & 35, 0 & 40, 1 & -22, 0 &  &  \\
b & 219, -129 & 240, -114 & 240, -114 & 136, -110 & 153, -114 & 139, -114 & 146, -114 & 218, -120 &  &  & 410, -114 & 294, -114 & 383, -105 & 265, -111 & 262, -112 & 247, -114 &  \\
b? adv & 224, -319 & 260, -110 & 217, -319 & 153, -110 & 147, -110 & 145, -110 & 145, -110 & 181, -110 & 103, -319 & 217, -319 & 382, -110 & 339, -319 & 339, -319 & 220, -319 & 220, -319 & 283, -110 & 238, -110 \\
b$\delta$ adv & \textbf{-5, 190} & \textbf{-20, -4} & \textbf{23, 205} & \textbf{-17, 0} & \textbf{6, -4} & -6, -4 & \textbf{1, -4} & 37, -10 &  &  & \textbf{28, -4} & \textbf{-45, 205} & \textbf{44, 214} & \textbf{45, 208} & 42, 207 & \textbf{-36, -4} &  \\
b? ink & 233, -319 & 264, -110 & 194, -319 & 164, -110 & 161, -110 & 145, -110 & 179, -110 & 193, -110 & 159, -319 & 215, -319 & 378, -110 & 347, -319 & 314, -319 & 190, -319 & 251, -319 & 286, -110 & 267, -110 \\
b$\delta$ ink & -14, 190 & -24, -4 & 46, 205 & -28, 0 & -8, -4 & -6, -4 & -33, -4 & \textbf{25, -10} &  &  & 32, -4 & -53, 205 & 69, 214 & 75, 208 & \textbf{11, 207} & -39, -4 &  \\

\end{tabular}
}
\end{table}
\clearpage


\begin{table}[p]
\captionsetup{justification=raggedright, singlelinecheck=false}
\caption{Lane comparison for Latin and IPA 6/8}
\label{table:lanes_latinIPA_6}
\setlength{\tabcolsep}{3.5pt}

{\footnotesize
\begin{tabular}{lrrrrrrrrrrrrrrrrr}
hex & 0276 & 0277 & 0278 & 0279 & 027A & 027B & 027C & 027D & 027E & 027F & 0280 & 0281 & 0282 & 0283 & 0284 & 0285 & 0286 \\
cat & Ll & Ll & Ll & Ll & Ll & Ll & Ll & Ll & Ll & Ll & Ll & Ll & Ll & Ll & Ll & Ll & Ll \\
\hline it\rule[-0.5em]{0pt}{1.9em} & {\itshape \char"0276 \char"0276\cgj\char"0307 \char"0276\cgj\char"0331} & {\itshape \char"0277 \char"0277\cgj\char"0307 \char"0277\cgj\char"0331} & {\itshape \char"0278 \char"0278\cgj\char"0307 \char"0278\cgj\char"0331} & {\itshape \char"0279 \char"0279\cgj\char"0307 \char"0279\cgj\char"0331} & {\itshape \char"027A \char"027A\cgj\char"0307 \char"027A\cgj\char"0331} & {\itshape \char"027B \char"027B\cgj\char"0307 \char"027B\cgj\char"0331} & {\itshape \char"027C \char"027C\cgj\char"0307 \char"027C\cgj\char"0331} & {\itshape \char"027D \char"027D\cgj\char"0307 \char"027D\cgj\char"0331} & {\itshape \char"027E \char"027E\cgj\char"0307 \char"027E\cgj\char"0331} & {\itshape \char"027F \char"027F\cgj\char"0307 \char"027F\cgj\char"0331} & {\itshape \char"0280 \char"0280\cgj\char"0307 \char"0280\cgj\char"0331} & {\itshape \char"0281 \char"0281\cgj\char"0307 \char"0281\cgj\char"0331} & {\itshape \char"0282 \char"0282\cgj\char"0307 \char"0282\cgj\char"0331} & {\itshape \char"0283 \char"0283\cgj\char"0307 \char"0283\cgj\char"0331} & {\itshape \char"0284 \char"0284\cgj\char"0307 \char"0284\cgj\char"0331} & {\itshape \char"0285 \char"0285\cgj\char"0307 \char"0285\cgj\char"0331} & {\itshape \char"0286 \char"0286\cgj\char"0307 \char"0286\cgj\char"0331} \\
patch\rule[-0.5em]{0pt}{1.9em} & {\LibertinusSerifPatch\itshape \char"0276 \char"0276\cgj\char"0307 \char"0276\cgj\char"0331} & {\LibertinusSerifPatch\itshape \char"0277 \char"0277\cgj\char"0307 \char"0277\cgj\char"0331} & {\LibertinusSerifPatch\itshape \char"0278 \char"0278\cgj\char"0307 \char"0278\cgj\char"0331} & {\LibertinusSerifPatch\itshape \char"0279 \char"0279\cgj\char"0307 \char"0279\cgj\char"0331} & {\LibertinusSerifPatch\itshape \char"027A \char"027A\cgj\char"0307 \char"027A\cgj\char"0331} & {\LibertinusSerifPatch\itshape \char"027B \char"027B\cgj\char"0307 \char"027B\cgj\char"0331} & {\LibertinusSerifPatch\itshape \char"027C \char"027C\cgj\char"0307 \char"027C\cgj\char"0331} & {\LibertinusSerifPatch\itshape \char"027D \char"027D\cgj\char"0307 \char"027D\cgj\char"0331} & {\LibertinusSerifPatch\itshape \char"027E \char"027E\cgj\char"0307 \char"027E\cgj\char"0331} & {\LibertinusSerifPatch\itshape \char"027F \char"027F\cgj\char"0307 \char"027F\cgj\char"0331} & {\LibertinusSerifPatch\itshape \char"0280 \char"0280\cgj\char"0307 \char"0280\cgj\char"0331} & {\LibertinusSerifPatch\itshape \char"0281 \char"0281\cgj\char"0307 \char"0281\cgj\char"0331} & {\LibertinusSerifPatch\itshape \char"0282 \char"0282\cgj\char"0307 \char"0282\cgj\char"0331} & {\LibertinusSerifPatch\itshape \char"0283 \char"0283\cgj\char"0307 \char"0283\cgj\char"0331} & {\LibertinusSerifPatch\itshape \char"0284 \char"0284\cgj\char"0307 \char"0284\cgj\char"0331} & {\LibertinusSerifPatch\itshape \char"0285 \char"0285\cgj\char"0307 \char"0285\cgj\char"0331} & {\LibertinusSerifPatch\itshape \char"0286 \char"0286\cgj\char"0307 \char"0286\cgj\char"0331} \\
a &  &  &  &  &  &  & 222, 703 &  &  &  &  &  &  &  &  &  &  \\
a? adv & 457, 645 & 494, 645 & 481, 890 & 315, 645 & 362, 819 & 352, 645 & 325, 645 & 320, 645 & 311, 645 & 307, 645 & 358, 645 & 349, 645 & 314, 645 & 343, 900 & 343, 900 & 343, 900 & 343, 900 \\
a$\delta$ adv &  &  &  &  &  &  & \textbf{-103, 58} &  &  &  &  &  &  &  &  &  &  \\
a? ink & 462, 645 & 497, 645 & 485, 890 & 223, 645 & 373, 819 & 362, 645 & 332, 645 & 331, 645 & 315, 645 & 311, 645 & 361, 645 & 351, 645 & 328, 645 & 336, 900 & 336, 900 & 335, 900 & 345, 900 \\
a$\delta$ ink &  &  &  &  &  &  & -110, 58 &  &  &  &  &  &  &  &  &  &  \\
b &  & 341, -114 &  & 190, -114 & 203, -114 &  & 127, -305 &  &  &  &  &  &  &  &  &  &  \\
b? adv & 296, -110 & 330, -126 & 231, -285 & 155, -110 & 161, -127 & 160, -259 & 128, -283 & 124, -279 & 151, -110 & 147, -110 & 197, -110 & 189, -110 & 121, -265 & 88, -299 & 88, -299 & 88, -299 & 92, -283 \\
b$\delta$ adv &  & 11, 12 &  & \textbf{35, -4} & 42, 13 &  & \textbf{-1, -22} &  &  &  &  &  &  &  &  &  &  \\
b? ink & 302, -110 & 333, -126 & 235, -285 & 63, -110 & 172, -127 & 170, -259 & 135, -283 & 135, -279 & 155, -110 & 151, -110 & 201, -110 & 191, -110 & 135, -265 & 81, -299 & 81, -299 & 80, -299 & 94, -283 \\
b$\delta$ ink &  & \textbf{8, 12} &  & 127, -4 & \textbf{31, 13} &  & -8, -22 &  &  &  &  &  &  &  &  &  &  \\
\hline si\rule[-0.5em]{0pt}{1.9em} & {\bfseries\itshape \char"0276 \char"0276\cgj\char"0307 \char"0276\cgj\char"0331} & {\bfseries\itshape \char"0277 \char"0277\cgj\char"0307 \char"0277\cgj\char"0331} & {\bfseries\itshape \char"0278 \char"0278\cgj\char"0307 \char"0278\cgj\char"0331} & {\bfseries\itshape \char"0279 \char"0279\cgj\char"0307 \char"0279\cgj\char"0331} & {\bfseries\itshape \char"027A \char"027A\cgj\char"0307 \char"027A\cgj\char"0331} & {\bfseries\itshape \char"027B \char"027B\cgj\char"0307 \char"027B\cgj\char"0331} & {\bfseries\itshape \char"027C \char"027C\cgj\char"0307 \char"027C\cgj\char"0331} & {\bfseries\itshape \char"027D \char"027D\cgj\char"0307 \char"027D\cgj\char"0331} & {\bfseries\itshape \char"027E \char"027E\cgj\char"0307 \char"027E\cgj\char"0331} & {\bfseries\itshape \char"027F \char"027F\cgj\char"0307 \char"027F\cgj\char"0331} & {\bfseries\itshape \char"0280 \char"0280\cgj\char"0307 \char"0280\cgj\char"0331} & {\bfseries\itshape \char"0281 \char"0281\cgj\char"0307 \char"0281\cgj\char"0331} & {\bfseries\itshape \char"0282 \char"0282\cgj\char"0307 \char"0282\cgj\char"0331} & {\bfseries\itshape \char"0283 \char"0283\cgj\char"0307 \char"0283\cgj\char"0331} & {\bfseries\itshape \char"0284 \char"0284\cgj\char"0307 \char"0284\cgj\char"0331} & {\bfseries\itshape \char"0285 \char"0285\cgj\char"0307 \char"0285\cgj\char"0331} & {\bfseries\itshape \char"0286 \char"0286\cgj\char"0307 \char"0286\cgj\char"0331} \\
patch\rule[-0.5em]{0pt}{1.9em} & {\LibertinusSerifPatch\bfseries\itshape \char"0276 \char"0276\cgj\char"0307 \char"0276\cgj\char"0331} & {\LibertinusSerifPatch\bfseries\itshape \char"0277 \char"0277\cgj\char"0307 \char"0277\cgj\char"0331} & {\LibertinusSerifPatch\bfseries\itshape \char"0278 \char"0278\cgj\char"0307 \char"0278\cgj\char"0331} & {\LibertinusSerifPatch\bfseries\itshape \char"0279 \char"0279\cgj\char"0307 \char"0279\cgj\char"0331} & {\LibertinusSerifPatch\bfseries\itshape \char"027A \char"027A\cgj\char"0307 \char"027A\cgj\char"0331} & {\LibertinusSerifPatch\bfseries\itshape \char"027B \char"027B\cgj\char"0307 \char"027B\cgj\char"0331} & {\LibertinusSerifPatch\bfseries\itshape \char"027C \char"027C\cgj\char"0307 \char"027C\cgj\char"0331} & {\LibertinusSerifPatch\bfseries\itshape \char"027D \char"027D\cgj\char"0307 \char"027D\cgj\char"0331} & {\LibertinusSerifPatch\bfseries\itshape \char"027E \char"027E\cgj\char"0307 \char"027E\cgj\char"0331} & {\LibertinusSerifPatch\bfseries\itshape \char"027F \char"027F\cgj\char"0307 \char"027F\cgj\char"0331} & {\LibertinusSerifPatch\bfseries\itshape \char"0280 \char"0280\cgj\char"0307 \char"0280\cgj\char"0331} & {\LibertinusSerifPatch\bfseries\itshape \char"0281 \char"0281\cgj\char"0307 \char"0281\cgj\char"0331} & {\LibertinusSerifPatch\bfseries\itshape \char"0282 \char"0282\cgj\char"0307 \char"0282\cgj\char"0331} & {\LibertinusSerifPatch\bfseries\itshape \char"0283 \char"0283\cgj\char"0307 \char"0283\cgj\char"0331} & {\LibertinusSerifPatch\bfseries\itshape \char"0284 \char"0284\cgj\char"0307 \char"0284\cgj\char"0331} & {\LibertinusSerifPatch\bfseries\itshape \char"0285 \char"0285\cgj\char"0307 \char"0285\cgj\char"0331} & {\LibertinusSerifPatch\bfseries\itshape \char"0286 \char"0286\cgj\char"0307 \char"0286\cgj\char"0331} \\
a &  &  &  &  &  &  & 222, 703 &  &  &  &  &  &  &  &  &  &  \\
a? adv & 313, 642 & 350, 645 & 328, 890 & 171, 640 & 217, 821 & 194, 645 & 181, 645 & 177, 645 & 167, 645 & 167, 645 & 214, 642 & 214, 645 & 178, 645 & 195, 890 & 195, 890 & 195, 890 & 189, 890 \\
a$\delta$ adv &  &  &  &  &  &  & 41, 58 &  &  &  &  &  &  &  &  &  &  \\
a? ink & 320, 642 & 356, 645 & 334, 890 & 158, 640 & 231, 821 & 218, 645 & 190, 645 & 190, 645 & 173, 645 & 169, 645 & 217, 642 & 210, 645 & 202, 645 & 190, 890 & 190, 890 & 190, 890 & 199, 890 \\
a$\delta$ ink &  &  &  &  &  &  & \textbf{32, 58} &  &  &  &  &  &  &  &  &  &  \\
b &  & 341, -114 &  & 190, -114 & 203, -114 &  & 127, -305 &  &  &  &  &  &  &  &  &  &  \\
b? adv & 322, -110 & 359, -110 & 251, -283 & 180, -110 & 190, -110 & 172, -257 & 155, -281 & 143, -319 & 176, -110 & 176, -110 & 223, -110 & 223, -110 & 155, -265 & 116, -296 & 116, -296 & 116, -296 & 105, -319 \\
b$\delta$ adv &  & \textbf{-18, -4} &  & \textbf{10, -4} & 13, -4 &  & \textbf{-28, -24} &  &  &  &  &  &  &  &  &  &  \\
b? ink & 329, -110 & 365, -110 & 257, -283 & 168, -110 & 204, -110 & 197, -257 & 164, -281 & 156, -319 & 182, -110 & 178, -110 & 226, -110 & 219, -110 & 179, -265 & 111, -296 & 111, -296 & 111, -296 & 115, -319 \\
b$\delta$ ink &  & -24, -4 &  & 22, -4 & \textbf{-1, -4} &  & -37, -24 &  &  &  &  &  &  &  &  &  &  \\

\end{tabular}
}
\end{table}
\clearpage


\begin{table}[p]
\captionsetup{justification=raggedright, singlelinecheck=false}
\caption{Lane comparison for Latin and IPA 7/8}
\label{table:lanes_latinIPA_7}
\setlength{\tabcolsep}{3.5pt}

{\footnotesize
\begin{tabular}{lrrrrrrrrrrrrrrrrr}
hex & 0287 & 0288 & 0289 & 028A & 028B & 028C & 028D & 028E & 028F & 0290 & 0291 & 0292 & 0293 & 0294 & 0295 & 0296 & 0297 \\
cat & Ll & Ll & Ll & Ll & Ll & Ll & Ll & Ll & Ll & Ll & Ll & Ll & Ll & Lu & Lu & Lu & Ll \\
\hline it\rule[-0.5em]{0pt}{1.9em} & {\itshape \char"0287 \char"0287\cgj\char"0307 \char"0287\cgj\char"0331} & {\itshape \char"0288 \char"0288\cgj\char"0307 \char"0288\cgj\char"0331} & {\itshape \char"0289 \char"0289\cgj\char"0307 \char"0289\cgj\char"0331} & {\itshape \char"028A \char"028A\cgj\char"0307 \char"028A\cgj\char"0331} & {\itshape \char"028B \char"028B\cgj\char"0307 \char"028B\cgj\char"0331} & {\itshape \char"028C \char"028C\cgj\char"0307 \char"028C\cgj\char"0331} & {\itshape \char"028D \char"028D\cgj\char"0307 \char"028D\cgj\char"0331} & {\itshape \char"028E \char"028E\cgj\char"0307 \char"028E\cgj\char"0331} & {\itshape \char"028F \char"028F\cgj\char"0307 \char"028F\cgj\char"0331} & {\itshape \char"0290 \char"0290\cgj\char"0307 \char"0290\cgj\char"0331} & {\itshape \char"0291 \char"0291\cgj\char"0307 \char"0291\cgj\char"0331} & {\itshape \char"0292 \char"0292\cgj\char"0307 \char"0292\cgj\char"0331} & {\itshape \char"0293 \char"0293\cgj\char"0307 \char"0293\cgj\char"0331} & {\itshape \char"0294 \char"0294\cgj\char"0307 \char"0294\cgj\char"0331} & {\itshape \char"0295 \char"0295\cgj\char"0307 \char"0295\cgj\char"0331} & {\itshape \char"0296 \char"0296\cgj\char"0307 \char"0296\cgj\char"0331} & {\itshape \char"0297 \char"0297\cgj\char"0307 \char"0297\cgj\char"0331} \\
patch\rule[-0.5em]{0pt}{1.9em} & {\LibertinusSerifPatch\itshape \char"0287 \char"0287\cgj\char"0307 \char"0287\cgj\char"0331} & {\LibertinusSerifPatch\itshape \char"0288 \char"0288\cgj\char"0307 \char"0288\cgj\char"0331} & {\LibertinusSerifPatch\itshape \char"0289 \char"0289\cgj\char"0307 \char"0289\cgj\char"0331} & {\LibertinusSerifPatch\itshape \char"028A \char"028A\cgj\char"0307 \char"028A\cgj\char"0331} & {\LibertinusSerifPatch\itshape \char"028B \char"028B\cgj\char"0307 \char"028B\cgj\char"0331} & {\LibertinusSerifPatch\itshape \char"028C \char"028C\cgj\char"0307 \char"028C\cgj\char"0331} & {\LibertinusSerifPatch\itshape \char"028D \char"028D\cgj\char"0307 \char"028D\cgj\char"0331} & {\LibertinusSerifPatch\itshape \char"028E \char"028E\cgj\char"0307 \char"028E\cgj\char"0331} & {\LibertinusSerifPatch\itshape \char"028F \char"028F\cgj\char"0307 \char"028F\cgj\char"0331} & {\LibertinusSerifPatch\itshape \char"0290 \char"0290\cgj\char"0307 \char"0290\cgj\char"0331} & {\LibertinusSerifPatch\itshape \char"0291 \char"0291\cgj\char"0307 \char"0291\cgj\char"0331} & {\LibertinusSerifPatch\itshape \char"0292 \char"0292\cgj\char"0307 \char"0292\cgj\char"0331} & {\LibertinusSerifPatch\itshape \char"0293 \char"0293\cgj\char"0307 \char"0293\cgj\char"0331} & {\LibertinusSerifPatch\itshape \char"0294 \char"0294\cgj\char"0307 \char"0294\cgj\char"0331} & {\LibertinusSerifPatch\itshape \char"0295 \char"0295\cgj\char"0307 \char"0295\cgj\char"0331} & {\LibertinusSerifPatch\itshape \char"0296 \char"0296\cgj\char"0307 \char"0296\cgj\char"0331} & {\LibertinusSerifPatch\itshape \char"0297 \char"0297\cgj\char"0307 \char"0297\cgj\char"0331} \\
a &  &  &  &  & 259, 645 &  &  &  &  &  &  & 328, 646 &  &  &  &  &  \\
a? adv & 288, 645 & 314, 776 & 398, 645 & 392, 645 & 382, 645 & 383, 645 & 506, 645 & 433, 838 & 372, 645 & 343, 645 & 352, 645 & 350, 645 & 355, 645 & 373, 850 & 373, 850 & 373, 850 & 411, 645 \\
a$\delta$ adv &  &  &  &  & \textbf{-123, 0} &  &  &  &  &  &  & -22, 1 &  &  &  &  &  \\
a? ink & 300, 645 & 337, 776 & 408, 645 & 396, 645 & 385, 645 & 383, 645 & 507, 645 & 430, 838 & 379, 645 & 391, 645 & 356, 645 & 347, 645 & 362, 645 & 384, 850 & 383, 850 & 376, 850 & 402, 645 \\
a$\delta$ ink &  &  &  &  & -126, 0 &  &  &  &  &  &  & \textbf{-19, 1} &  &  &  &  &  \\
b &  &  &  & 258, -114 & 251, -114 &  &  &  & 239, -107 &  &  & 130, -319 &  & 202, -107 & 230, -107 &  &  \\
b? adv & 100, -239 & 83, -311 & 238, -110 & 229, -121 & 218, -124 & 222, -110 & 346, -110 & 231, -110 & 212, -110 & 147, -276 & 171, -207 & 145, -319 & 133, -397 & 169, -110 & 169, -110 & 164, -131 & 206, -319 \\
b$\delta$ adv &  &  &  & 29, 7 & 33, 10 &  &  &  & 27, 3 &  &  & -15, 0 &  & 33, 3 & 61, 3 &  &  \\
b? ink & 112, -239 & 106, -311 & 248, -110 & 233, -121 & 222, -124 & 223, -110 & 347, -110 & 229, -110 & 219, -110 & 195, -276 & 175, -207 & 142, -319 & 141, -397 & 180, -110 & 179, -110 & 167, -131 & 197, -319 \\
b$\delta$ ink &  &  &  & \textbf{25, 7} & \textbf{29, 10} &  &  &  & \textbf{20, 3} &  &  & \textbf{-12, 0} &  & \textbf{22, 3} & \textbf{51, 3} &  &  \\
\hline si\rule[-0.5em]{0pt}{1.9em} & {\bfseries\itshape \char"0287 \char"0287\cgj\char"0307 \char"0287\cgj\char"0331} & {\bfseries\itshape \char"0288 \char"0288\cgj\char"0307 \char"0288\cgj\char"0331} & {\bfseries\itshape \char"0289 \char"0289\cgj\char"0307 \char"0289\cgj\char"0331} & {\bfseries\itshape \char"028A \char"028A\cgj\char"0307 \char"028A\cgj\char"0331} & {\bfseries\itshape \char"028B \char"028B\cgj\char"0307 \char"028B\cgj\char"0331} & {\bfseries\itshape \char"028C \char"028C\cgj\char"0307 \char"028C\cgj\char"0331} & {\bfseries\itshape \char"028D \char"028D\cgj\char"0307 \char"028D\cgj\char"0331} & {\bfseries\itshape \char"028E \char"028E\cgj\char"0307 \char"028E\cgj\char"0331} & {\bfseries\itshape \char"028F \char"028F\cgj\char"0307 \char"028F\cgj\char"0331} & {\bfseries\itshape \char"0290 \char"0290\cgj\char"0307 \char"0290\cgj\char"0331} & {\bfseries\itshape \char"0291 \char"0291\cgj\char"0307 \char"0291\cgj\char"0331} & {\bfseries\itshape \char"0292 \char"0292\cgj\char"0307 \char"0292\cgj\char"0331} & {\bfseries\itshape \char"0293 \char"0293\cgj\char"0307 \char"0293\cgj\char"0331} & {\bfseries\itshape \char"0294 \char"0294\cgj\char"0307 \char"0294\cgj\char"0331} & {\bfseries\itshape \char"0295 \char"0295\cgj\char"0307 \char"0295\cgj\char"0331} & {\bfseries\itshape \char"0296 \char"0296\cgj\char"0307 \char"0296\cgj\char"0331} & {\bfseries\itshape \char"0297 \char"0297\cgj\char"0307 \char"0297\cgj\char"0331} \\
patch\rule[-0.5em]{0pt}{1.9em} & {\LibertinusSerifPatch\bfseries\itshape \char"0287 \char"0287\cgj\char"0307 \char"0287\cgj\char"0331} & {\LibertinusSerifPatch\bfseries\itshape \char"0288 \char"0288\cgj\char"0307 \char"0288\cgj\char"0331} & {\LibertinusSerifPatch\bfseries\itshape \char"0289 \char"0289\cgj\char"0307 \char"0289\cgj\char"0331} & {\LibertinusSerifPatch\bfseries\itshape \char"028A \char"028A\cgj\char"0307 \char"028A\cgj\char"0331} & {\LibertinusSerifPatch\bfseries\itshape \char"028B \char"028B\cgj\char"0307 \char"028B\cgj\char"0331} & {\LibertinusSerifPatch\bfseries\itshape \char"028C \char"028C\cgj\char"0307 \char"028C\cgj\char"0331} & {\LibertinusSerifPatch\bfseries\itshape \char"028D \char"028D\cgj\char"0307 \char"028D\cgj\char"0331} & {\LibertinusSerifPatch\bfseries\itshape \char"028E \char"028E\cgj\char"0307 \char"028E\cgj\char"0331} & {\LibertinusSerifPatch\bfseries\itshape \char"028F \char"028F\cgj\char"0307 \char"028F\cgj\char"0331} & {\LibertinusSerifPatch\bfseries\itshape \char"0290 \char"0290\cgj\char"0307 \char"0290\cgj\char"0331} & {\LibertinusSerifPatch\bfseries\itshape \char"0291 \char"0291\cgj\char"0307 \char"0291\cgj\char"0331} & {\LibertinusSerifPatch\bfseries\itshape \char"0292 \char"0292\cgj\char"0307 \char"0292\cgj\char"0331} & {\LibertinusSerifPatch\bfseries\itshape \char"0293 \char"0293\cgj\char"0307 \char"0293\cgj\char"0331} & {\LibertinusSerifPatch\bfseries\itshape \char"0294 \char"0294\cgj\char"0307 \char"0294\cgj\char"0331} & {\LibertinusSerifPatch\bfseries\itshape \char"0295 \char"0295\cgj\char"0307 \char"0295\cgj\char"0331} & {\LibertinusSerifPatch\bfseries\itshape \char"0296 \char"0296\cgj\char"0307 \char"0296\cgj\char"0331} & {\LibertinusSerifPatch\bfseries\itshape \char"0297 \char"0297\cgj\char"0307 \char"0297\cgj\char"0331} \\
a &  &  &  &  & 259, 645 &  &  &  &  &  &  & 218, 646 &  &  &  &  &  \\
a? adv & 149, 645 & 172, 787 & 255, 645 & 248, 642 & 241, 645 & 239, 645 & 363, 645 & 286, 834 & 228, 645 & 202, 645 & 208, 645 & 206, 645 & 211, 645 & 227, 850 & 227, 850 & 227, 850 & 268, 645 \\
a$\delta$ adv &  &  &  &  & 18, 0 &  &  &  &  &  &  & \textbf{12, 1} &  &  &  &  &  \\
a? ink & 136, 645 & 195, 787 & 285, 645 & 255, 642 & 246, 645 & 220, 645 & 346, 645 & 273, 834 & 237, 645 & 252, 645 & 213, 645 & 205, 645 & 217, 645 & 244, 850 & 243, 850 & 232, 850 & 258, 645 \\
a$\delta$ ink &  &  &  &  & \textbf{13, 0} &  &  &  &  &  &  & 13, 1 &  &  &  &  &  \\
b &  &  &  & 258, -114 & 251, -114 &  &  &  & 239, -107 &  &  & 141, -319 &  & 202, -107 & 230, -107 &  &  \\
b? adv & 129, -250 & 109, -319 & 264, -110 & 257, -110 & 249, -110 & 247, -110 & 371, -110 & 256, -110 & 237, -110 & 176, -279 & 196, -210 & 172, -319 & 157, -413 & 194, -110 & 194, -110 & 194, -110 & 237, -302 \\
b$\delta$ adv &  &  &  & \textbf{1, -4} & \textbf{2, -4} &  &  &  & \textbf{2, 3} &  &  & -31, 0 &  & \textbf{8, 3} & 36, 3 &  &  \\
b? ink & 116, -250 & 132, -319 & 294, -110 & 264, -110 & 255, -110 & 229, -110 & 355, -110 & 243, -110 & 246, -110 & 226, -279 & 201, -210 & 171, -319 & 164, -413 & 211, -110 & 210, -110 & 199, -110 & 228, -302 \\
b$\delta$ ink &  &  &  & -6, -4 & -4, -4 &  &  &  & -7, 3 &  &  & \textbf{-30, 0} &  & -9, 3 & \textbf{20, 3} &  &  \\

\end{tabular}
}
\end{table}
\clearpage


\begin{table}[p]
\captionsetup{justification=raggedright, singlelinecheck=false}
\caption{Lane comparison for Latin and IPA 8/8}
\label{table:lanes_latinIPA_8}
\setlength{\tabcolsep}{3.5pt}

{\footnotesize
\begin{tabular}{lrrrrrrrrrrrrrr}
hex & 0299 & 029B & 029C & 029D & 029F & 02AE & 02AF & 03B2 & 03B8 & 03C7 & 1EFA & 2132 & 2C69 & 2C6B \\
cat & Ll & Ll & Ll & Ll & Ll & Ll & Ll & Ll & Ll & Ll & Lu & Lu & Lu & Lu \\
\hline it\rule[-0.5em]{0pt}{1.9em} & {\itshape \char"0299 \char"0299\cgj\char"0307 \char"0299\cgj\char"0331} & {\itshape \char"029B \char"029B\cgj\char"0307 \char"029B\cgj\char"0331} & {\itshape \char"029C \char"029C\cgj\char"0307 \char"029C\cgj\char"0331} & {\itshape \char"029D \char"029D\cgj\char"0307 \char"029D\cgj\char"0331} & {\itshape \char"029F \char"029F\cgj\char"0307 \char"029F\cgj\char"0331} & {\itshape \char"02AE \char"02AE\cgj\char"0307 \char"02AE\cgj\char"0331} & {\itshape \char"02AF \char"02AF\cgj\char"0307 \char"02AF\cgj\char"0331} & {\itshape \char"03B2 \char"03B2\cgj\char"0307 \char"03B2\cgj\char"0331} & {\itshape \char"03B8 \char"03B8\cgj\char"0307 \char"03B8\cgj\char"0331} & {\itshape \char"03C7 \char"03C7\cgj\char"0307 \char"03C7\cgj\char"0331} & {\itshape \char"1EFA \char"1EFA\cgj\char"0307 \char"1EFA\cgj\char"0331} & {\itshape \char"2132 \char"2132\cgj\char"0307 \char"2132\cgj\char"0331} & {\itshape \char"2C69 \char"2C69\cgj\char"0307 \char"2C69\cgj\char"0331} & {\itshape \char"2C6B \char"2C6B\cgj\char"0307 \char"2C6B\cgj\char"0331} \\
patch\rule[-0.5em]{0pt}{1.9em} & {\LibertinusSerifPatch\itshape \char"0299 \char"0299\cgj\char"0307 \char"0299\cgj\char"0331} & {\LibertinusSerifPatch\itshape \char"029B \char"029B\cgj\char"0307 \char"029B\cgj\char"0331} & {\LibertinusSerifPatch\itshape \char"029C \char"029C\cgj\char"0307 \char"029C\cgj\char"0331} & {\LibertinusSerifPatch\itshape \char"029D \char"029D\cgj\char"0307 \char"029D\cgj\char"0331} & {\LibertinusSerifPatch\itshape \char"029F \char"029F\cgj\char"0307 \char"029F\cgj\char"0331} & {\LibertinusSerifPatch\itshape \char"02AE \char"02AE\cgj\char"0307 \char"02AE\cgj\char"0331} & {\LibertinusSerifPatch\itshape \char"02AF \char"02AF\cgj\char"0307 \char"02AF\cgj\char"0331} & {\LibertinusSerifPatch\itshape \char"03B2 \char"03B2\cgj\char"0307 \char"03B2\cgj\char"0331} & {\LibertinusSerifPatch\itshape \char"03B8 \char"03B8\cgj\char"0307 \char"03B8\cgj\char"0331} & {\LibertinusSerifPatch\itshape \char"03C7 \char"03C7\cgj\char"0307 \char"03C7\cgj\char"0331} & {\LibertinusSerifPatch\itshape \char"1EFA \char"1EFA\cgj\char"0307 \char"1EFA\cgj\char"0331} & {\LibertinusSerifPatch\itshape \char"2132 \char"2132\cgj\char"0307 \char"2132\cgj\char"0331} & {\LibertinusSerifPatch\itshape \char"2C69 \char"2C69\cgj\char"0307 \char"2C69\cgj\char"0331} & {\LibertinusSerifPatch\itshape \char"2C6B \char"2C6B\cgj\char"0307 \char"2C6B\cgj\char"0331} \\
a &  &  &  &  &  & 278, -214 & 278, -214 &  &  &  &  &  & 336, 798 & 320, 803 \\
a? adv & 352, 645 & 410, 796 & 416, 645 & 303, 810 & 355, 645 & 422, 645 & 433, 645 & 425, 890 & 413, 890 & 372, 645 &  &  & 516, 850 & 490, 850 \\
a$\delta$ adv &  &  &  &  &  & \textbf{-144, -859} & \textbf{-155, -859} &  &  &  &  &  & \textbf{-180, -52} & \textbf{-170, -47} \\
a? ink & 354, 645 & 450, 796 & 422, 645 & 292, 810 & 345, 645 & 432, 645 & 467, 645 & 428, 890 & 419, 890 & 375, 645 &  &  & 521, 850 & 505, 850 \\
a$\delta$ ink &  &  &  &  &  & -154, -859 & -189, -859 &  &  &  &  &  & -185, -52 & -185, -47 \\
b &  &  & 273, -110 &  &  & 318, 542 & 318, 542 & 275, -114 &  &  &  &  & 342, -109 & 287, -107 \\
b? adv & 192, -110 & 217, -110 & 256, -110 & 77, -254 & 194, -110 & 220, -305 & 227, -326 & 168, -319 & 200, -110 & 167, -319 &  &  & 288, -223 & 265, -207 \\
b$\delta$ adv &  &  & 17, 0 &  &  & 98, 847 & 91, 868 & 107, 205 &  &  &  &  & 54, 114 & 22, 100 \\
b? ink & 194, -110 & 258, -110 & 262, -110 & 66, -254 & 185, -110 & 230, -305 & 261, -326 & 171, -319 & 207, -110 & 170, -319 &  &  & 293, -223 & 280, -207 \\
b$\delta$ ink &  &  & \textbf{11, 0} &  &  & \textbf{88, 847} & \textbf{57, 868} & \textbf{104, 205} &  &  &  &  & \textbf{49, 114} & \textbf{7, 100} \\
\hline si\rule[-0.5em]{0pt}{1.9em} & {\bfseries\itshape \char"0299 \char"0299\cgj\char"0307 \char"0299\cgj\char"0331} & {\bfseries\itshape \char"029B \char"029B\cgj\char"0307 \char"029B\cgj\char"0331} & {\bfseries\itshape \char"029C \char"029C\cgj\char"0307 \char"029C\cgj\char"0331} & {\bfseries\itshape \char"029D \char"029D\cgj\char"0307 \char"029D\cgj\char"0331} & {\bfseries\itshape \char"029F \char"029F\cgj\char"0307 \char"029F\cgj\char"0331} & {\bfseries\itshape \char"02AE \char"02AE\cgj\char"0307 \char"02AE\cgj\char"0331} & {\bfseries\itshape \char"02AF \char"02AF\cgj\char"0307 \char"02AF\cgj\char"0331} & {\bfseries\itshape \char"03B2 \char"03B2\cgj\char"0307 \char"03B2\cgj\char"0331} & {\bfseries\itshape \char"03B8 \char"03B8\cgj\char"0307 \char"03B8\cgj\char"0331} & {\bfseries\itshape \char"03C7 \char"03C7\cgj\char"0307 \char"03C7\cgj\char"0331} & {\bfseries\itshape \char"1EFA \char"1EFA\cgj\char"0307 \char"1EFA\cgj\char"0331} & {\bfseries\itshape \char"2132 \char"2132\cgj\char"0307 \char"2132\cgj\char"0331} & {\bfseries\itshape \char"2C69 \char"2C69\cgj\char"0307 \char"2C69\cgj\char"0331} & {\bfseries\itshape \char"2C6B \char"2C6B\cgj\char"0307 \char"2C6B\cgj\char"0331} \\
patch\rule[-0.5em]{0pt}{1.9em} & {\LibertinusSerifPatch\bfseries\itshape \char"0299 \char"0299\cgj\char"0307 \char"0299\cgj\char"0331} & {\LibertinusSerifPatch\bfseries\itshape \char"029B \char"029B\cgj\char"0307 \char"029B\cgj\char"0331} & {\LibertinusSerifPatch\bfseries\itshape \char"029C \char"029C\cgj\char"0307 \char"029C\cgj\char"0331} & {\LibertinusSerifPatch\bfseries\itshape \char"029D \char"029D\cgj\char"0307 \char"029D\cgj\char"0331} & {\LibertinusSerifPatch\bfseries\itshape \char"029F \char"029F\cgj\char"0307 \char"029F\cgj\char"0331} & {\LibertinusSerifPatch\bfseries\itshape \char"02AE \char"02AE\cgj\char"0307 \char"02AE\cgj\char"0331} & {\LibertinusSerifPatch\bfseries\itshape \char"02AF \char"02AF\cgj\char"0307 \char"02AF\cgj\char"0331} & {\LibertinusSerifPatch\bfseries\itshape \char"03B2 \char"03B2\cgj\char"0307 \char"03B2\cgj\char"0331} & {\LibertinusSerifPatch\bfseries\itshape \char"03B8 \char"03B8\cgj\char"0307 \char"03B8\cgj\char"0331} & {\LibertinusSerifPatch\bfseries\itshape \char"03C7 \char"03C7\cgj\char"0307 \char"03C7\cgj\char"0331} & {\LibertinusSerifPatch\bfseries\itshape \char"1EFA \char"1EFA\cgj\char"0307 \char"1EFA\cgj\char"0331} & {\LibertinusSerifPatch\bfseries\itshape \char"2132 \char"2132\cgj\char"0307 \char"2132\cgj\char"0331} & {\LibertinusSerifPatch\bfseries\itshape \char"2C69 \char"2C69\cgj\char"0307 \char"2C69\cgj\char"0331} & {\LibertinusSerifPatch\bfseries\itshape \char"2C6B \char"2C6B\cgj\char"0307 \char"2C6B\cgj\char"0331} \\
a &  &  &  &  &  & 278, -214 & 278, -214 &  &  &  &  &  & 336, 798 & 320, 803 \\
a? adv & 213, 645 & 278, 796 & 342, 671 & 158, 812 & 211, 645 & 279, 645 & 289, 645 & 302, 890 & 277, 873 & 256, 645 &  &  & 358, 850 & 332, 850 \\
a$\delta$ adv &  &  &  &  &  & \textbf{-1, -859} & \textbf{-11, -859} &  &  &  &  &  & \textbf{-22, -52} & \textbf{-12, -47} \\
a? ink & 220, 645 & 325, 796 & 344, 671 & 147, 812 & 202, 645 & 292, 645 & 324, 645 & 321, 890 & 286, 873 & 272, 645 &  &  & 361, 850 & 345, 850 \\
a$\delta$ ink &  &  &  &  &  & -14, -859 & -46, -859 &  &  &  &  &  & -25, -52 & -25, -47 \\
b &  &  & 327, -110 &  &  & 318, 542 & 318, 542 & 275, -114 &  &  &  &  & 342, -109 & 287, -107 \\
b? adv & 221, -110 & 256, -110 & 345, -110 & 103, -252 & 219, -110 & 245, -319 & 255, -319 & 218, -319 & 239, -110 & 231, -275 &  &  & 302, -221 & 279, -206 \\
b$\delta$ adv &  &  & \textbf{-18, 0} &  &  & 73, 861 & 63, 861 & 57, 205 &  &  &  &  & 40, 112 & 8, 99 \\
b? ink & 229, -110 & 303, -110 & 347, -110 & 93, -252 & 211, -110 & 258, -319 & 290, -319 & 237, -319 & 248, -110 & 247, -275 &  &  & 305, -221 & 292, -206 \\
b$\delta$ ink &  &  & -20, 0 &  &  & \textbf{60, 861} & \textbf{28, 861} & \textbf{38, 205} &  &  &  &  & \textbf{37, 112} & \textbf{-5, 99} \\

\end{tabular}
}
\end{table}
\clearpage


\end{document}
